\documentclass[../main.tex]{subfiles}
\begin{document}
%==========================================TIÊU ĐỀ TẬP========================================
\begin{table}[h]
  \begin{adjustwidth}{-1cm}{}
    \begin{tabular}{>{\centering\arraybackslash}p{6cm}|>{\raggedright\arraybackslash}p{14cm}}
      \multirow{5}{6cm}
      % Cột bên trái: ÉPISODE và số tập
      {\\[-40pt]
        \flaregothic\fontsize{20pt}{20pt}\selectfont \centering
        \color{eptype}HISTOIRE\color{black}\\
        \gangofthree\fontsize{72pt}{72pt}\selectfont
        \color{epnum}5
        \\[18pt]
        % \flaregothic\fontsize{14pt}{14pt}\selectfont
        % \color{epattr}BIENVENUE, \uline{\color{Banpire} Mel} !
        % \flaregothic\fontsize{18pt}{18pt}\selectfont
        % \color{epattr}ÉPISODE PERDU
        \color{black}
      }
       &
      % Dòng thứ nhất: tên tập tiếng Nhật
      {\fotskip\fontsize{18pt}{18pt}\selectfont \ruby{元}{がん}\ruby{日}{じつ}の\ruby{朝}{あさ}、\ruby{父}{ちち}に\ruby{会}{あ}う}
      \\[6pt]
      % Dòng thứ hai: tên tập tiếng Anh
       & {\cafeteria\fontsize{18pt}{18pt}\selectfont New Year's Day with Father!}         \\[4pt]
      % Dòng thứ ba: tên tập tiếng Việt
       & {\vnmsans\fontsize{18pt}{18pt}\selectfont Mùng Một Tết Cha!}                     \\[8pt]
      % Dòng thứ tư: Nhân vật xuất hiện
       & {\brian{\fontsize{14pt}{14pt}\selectfont Track used: Lunar Zoo Year - Loonboon}}
    \end{tabular}
  \end{adjustwidth}
\end{table}
%===========================================NỘI DUNG==========================================
\color{black}

\pov{Hikawa Kuon}

Sau ``tai nạn pháo hoa'' vừa rồi, ngoại trừ Pekora-chan đã ngất xỉu (hiện đang nằm thẳng
cẳng trong phòng tôi) và Haato-chan với Shion-chan dưới nhà canh nồi bánh chưng, đa
số các thành viên khác đã về ``doanh trại'' để nghỉ ngơi.

Tôi cá là bây giờ, chỉ cần ghé tai vào cửa, kiểu gì cũng nghe được ``dàn giao hưởng'' ngáy
đầy ``du dương'' của A-san dành cho các cô gái của công ty nhà ta. Chằng biết họ có hối
lỗi sau vụ việc nghiêm trọng này không, nhưng dám chắc họ vẫn sẽ ``báo'' chúng tôi dài
dài\dots

Và khi tôi nói ``đa số,'' tức là vẫn còn vài người kiên trì ở lại. Đó là Mel-chan, Mio-chan,
và Sora-chan, ba cô gái mà nếu phải đặt biệt danh, tôi sẽ gọi là ``học sinh gương mẫu''
trong nhóm.

Cùng gia đình tôi, ba cô gái này xuống tầng bốn để chào hỏi ông bà nội chúng tôi và các
bác đang quây quần dưới nhà. Ngay khi vừa bước vào phòng khách, hương thơm của trà
mạn nóng hòa quyện với mùi bánh mứt và không khí ấm cúng của ngày Tết ùa tới.

Chúng tôi lần lượt cúi chào mọi người. Ngoài bác Hijaki-san - người anh thứ trong dòng họ Hikawa (nếu tính ông nội chúng tôi) - ngoài ra còn có bác Hanasaki (\ruby{花}{はな}\ruby{さき}{崎}), vợ bác ấy, cùng hai
cô con gái của họ: Hokokata-san (\ruby{埃}{ほこ}\ruby{方}{かた}, chị cả) và Busako-san (\ruby{部}{ぶ}\ruby{迫}{さこ}, cô em).

Mọi thứ diễn ra theo đúng nhịp điệu quen thuộc của những cái Tết xưa nay, ngoại trừ việc
ông nội tôi giờ đã không còn nói chuyện. Có lẽ vì tuổi cao sức yếu, ông chỉ dùng cái gật
đầu hoặc lắc đầu để trả lời câu hỏi của con cháu. Tuy vậy, ánh mắt ông vẫn ánh lên nét
vui vẻ mỗi khi nhìn thấy các cháu chắt quây quần.

Vì Hokoru-san là người duy nhất ngoài gia đình tôi (và bác Hijaki-san) biết tiếng Nhật,
nên cả tôi và chị ấy phải kiêm thêm vai trò thông dịch viên bất đắc dĩ. Nếu các bạn đang
tự hỏi: ``Ơ, nhưng mà, tôi tưởng các cô gái có thể giao tiếp bằng tiếng bản xứ tại đây do
đã ăn bánh mì chuyển ngữ chứ?'' thì rất tiếc, món bảo bối đó chỉ có hiệu lực trong mười
hai giờ, và tôi chưa kịp làm thêm, nên tôi mới phải nhờ chị ấy làm hộ công việc này.

Chúng tôi giới thiệu với mọi người từng người một: Mel-chan với nụ cười tươi như nắng,
Mio-chan có phần rụt rè, và Sora-chan với thái độ lễ phép đáng khen. Sau đó, câu chuyện
chuyển sang những chủ đề xã giao nhẹ nhàng: công việc, đời sống, và vài câu đùa dịp Tết
đặc trưng ở xứ Etsunan.

Mọi chuyện đang êm đềm thì bỗng nhiên, như muốn khuấy động không khí, Busako-san
hỏi một câu khiến ai nấy đều khựng lại: ``{\color{red} Thế các em đã có gấu chưa?}''

Cả phòng bật cười, trừ ba cô gái. Họ ngẩn ra, nhìn nhau rồi quay sang tôi, rõ ràng không
hiểu chị ấy vừa nói gì. Tôi đành phải giải thích: ``Chị ấy đang hỏi xem mấy đứa đã có bạn
trai chưa đó.''

Nghe tôi nói, mặt cả ba lập tức đỏ bừng như vừa bị ai trêu. Mel-chan nở nụ cười ngượng
ngùng, Mio-chan đứng đơ như tượng, còn Sora-chan thì lắp bắp trả lời.

\begin{dialogue}
  \speak{Sora} E-Em chưa có ạ\dots
  \speak{Mio} *đỏ mặt* \dots
  \speak{Mel} Em cũng chưa ạ\dots *cười trừ*
\end{dialogue}

Tôi không hề thấy ngạc nhiên. Là những thần tượng, họ bị ràng buộc bởi hình tượng ``của chung.''
Điều này chẳng phải điều mới mẻ gì, nhưng để giải thích cho các bác hiểu, tôi đành phải
lên tiếng: ``{\color{red} Họ là idol ạ, nên quy tắc công việc không cho phép họ hẹn hò.}''

Bác Hanasaki-san nghe vậy cũng gật gù hiểu, rồi cất giọng hỏi: ``{\color{red} Vậy đến khi nào các cháu mới tìm được nửa kia?}''

Tôi đáp, nở một nụ cười trừ: ``{\color{red}Cháu nghĩ chắc phải đợi họ giải nghệ. Nhưng mà nhìn thế này thì chắc còn lâu lắm ạ.}''

Hai người phụ nữ nghe tôi giải đáp thế, cười xòa với nhau. Họ bắt đầu bàn đến chuyện
yêu đương, một trong những chủ đề mà tôi và các cô gái đếu chưa hề để tâm đến\dots

Tôi quay sang ba người họ, đang tỏ vẻ bối rối hơn bao giờ hết, và nói như để trấn an:
``Đấy, các em thấy chưa, người ta hay hỏi những cái bất lịch sự như vậy đấy.''

Mio-chan bày tỏ bức xúc: ``Nhưng chẳng phải như thế là bất lịch sự sao?'' Tôi chỉ có thể
thở dài đáp: ``Vì thế người ta mới càng hỏi. Ở quê anh gần như ai cũng vậy hết mà.''

Mel-chan tỏ vẻ ngạc nhiên: ``Anh trông có vẻ quen với mấy câu hỏi này nhỉ?'' Tôi chẹp
miệng, nhìn mẹ tôi và bác Hasanaki-san ``buôn dưa lê,'' rồi thở dài: ``Hồi trước anh cũng
toàn bị hỏi có người yêu chưa. Hiểu cảm giác này lắm mà. Sau này có khi họ còn hỏi có
lương chưa nữa kìa\dots''

Và giờ đến ``tiết mục'' mà đám trẻ con (bao gồm tôi trước kia) mong chờ nhất: nhận bao
lì xì. Không khí ấm cúng của buổi họp mặt bỗng trở nên háo hức hơn khi bà nội tôi (tôi sẽ
gọi tắt là nội) chậm rãi rút ra một xấp tiền năm mươi nghìn đồng, nói với giọng đầy hóm
hỉnh: ``{\color{red} Nhân dịp năm mới, tao đây có tiền mừng tuổi cho chúng mày này.}''

Bà lần lượt chia tiền cho tất cả mọi người trong nhà, không phân biệt ai là bác, bố mẹ hay
con cháu. Đương nhiên, Mio-chan, Mel-chan, và Sora-chan cũng không bị bỏ sót.

Sora-chan, với nét mặt ngạc nhiên lẫn tò mò, hỏi tôi: ``Ở đây bọn anh cũng lì xì ạ?''

Tôi gật đầu xác nhận, rồi em ấy tiếp tục hỏi với vẻ bối rối: ``Nhưng bình thường họ sẽ cho
vào một cái bao, rồi mới đưa cho mọi người mà?''

Tôi khẽ cười giải thích: ``Ở nhà anh thì là như thế này. Còn bình thường, như em nói,
người ta sẽ cho tiền vào một cái bao màu đỏ trước khi lì xì.''

Mel-chan nghiêng đầu hỏi thêm: ``Màu đỏ tượng trưng cho may mắn, đúng không nhỉ?''
Tôi đáp: ``Chính xác. Ở Nhật, bọn em gọi tiền này là \uline{otoshidama} (お年玉), còn ở đây bọn
anh gọi là \uline{tiền mừng tuổi}.''

Mio-chan lặp lại cụm từ tôi vừa nói đó, tỏ vẻ tò mò: ``Mừng tuổi?'' Cô em gái tôi tiếp lời:
``Tức là chúc mừng tuổi mới đấy các chị. Hiểu nôm na tức là lộc đầu năm đấy ạ.''

Cả ba người gật gù, tờ mờ hiểu hơn với ý nghĩa của tiền lì xì ở xứ Etsunan. ``Nghe có vẻ
giống phong tục ở Nhật nhỉ,'' nàng Thần tượng cất lời.

``Thì, bởi vì Etsunan với Nhật Bản đã từng-'' tôi định giải thích thêm, thì bỗng bà nội cất
tiếng, làm chúng tôi giật mình: ``{\color{red} Chúng mày có cần thêm không?}''

Lần này, các bác, các chị, bố mẹ, và cả hai anh em tôi đồng thanh phản ứng, gần như bản
năng: ``{\color{red} Dạ thôi, không cần đâu ạ!}''

Bà nội ngơ ngác, tay khẽ lật lại xấp tiền: ``{\color{red} Ủa, tưởng tao chưa cho mà\dots}''

Cả tôi và ba cô gái cùng đồng loạt thốt lên: ``\textbf{Bà cho rồi ạ!}'' Bà khựng lại, rồi gật đầu
cười gượng: ``À, thế thôi vậy.''

Tôi quay sang nhìn Sói đen, cười nhẹ và tiếp tục hỏi về chuyện lì xì: ``Em có vẻ cũng tìm
hiểu nhiều vè ý nghĩa lì xì ngày Tết nhỉ?''

Em đáp, vẻ mặt vẫn giữ nét điềm tĩnh thường thấy: ``Em đã ăn Tết hai năm ở Nhật rồi,
nên em cũng hiểu được một phần ý nghĩa của nó.''

Mel-chan mỉm cười xen vào: ``Với lại em nghĩ cái bao mới là quan trọng nhất. Tiền bên
trong bao nhiêu cũng không phải điều cần thiết.''

Tôi gật đầu đồng tình, rồi thở dài như đang hồi tưởng: ``Đúng thế, quan trọng là cái tấm
lòng của người ta. Nhưng không giống bọn trẻ con bây giờ\dots''

Thấy tôi đang suy tư, thi thoảng còn lắc đầu khe khẽ, Sora-chan tò mò hỏi: ``Ý anh là sao, Kuon-senpai?''

``Bây giờ cứ nhắc đến lì xì là đủ vấn đề\dots\ Phong tục đã bị bóp méo và biến chất so với ý
nghĩa ban đầu. Bây giờ cứ mỗi lần bố mẹ lì xì thì câu cửa miệng người ta toàn hỏi là `lì
xì được bao nhiêu.' rồi còn tính toán xem lì xì như thế nào để được lợi nữa.''

``Nói tóm lại là, hỏng hết cả một thế hệ ạ,'' cô em gái tôi kết luận.

Sora-chan khẽ nhăn mặt đượm buồn, dường như đồng cảm: ``Anh nói cũng phải\dots\ Em
nghe nói hồi trước Shion-chan, Subaru-chan, Ayame-chan với Miko-chan chơi xổ số\footnote{Xem {\flaregothic ÉPISODE 4}, quyển số Không.}.
Họ còn tranh nhau vé ác lắm.''

Mio-chan liếc nhìn tôi, thở dài ngao ngán: ``Đến mức độ làm sập cả công ty là đủ hiểu
rồi.'' Dù tôi không có mặt ở đó, nhưng nhớ lại cách mà mọi người kể với tôi ngày hôm
đó cũng đủ để khiến tôi ngán ngẩm mà tặc lưỡi: ``\dots Chán.''

Để đáp lại lòng thành của bà nội và các bác, ba cô gái bất ngờ lấy ra những phong bao
lì xì được chuẩn bị sẵn. Lần này, chúng tôi lì xì lại họ bằng tiền đồng với mệnh giá cao
hơn, khiến cả nhà không khỏi ngạc nhiên. Tôi không biết họ đổi tiền từ lúc nào, nhưng rõ
ràng, sự chu đáo này khiến mọi người đều cảm kích.

Chúng tôi tiếp tục trò chuyện rôm rả đến tận một giờ sáng. Khi này, mẹ tôi đã xuống tầng
một cùng bà ngoại để trông nồi bánh chưng với Phù thủy-chan và Haachama-chan, còn
lại cả nhóm và ba bố con chúng tôi lên phòng nghỉ ngơi, chuẩn bị cho buổi sáng đầu năm.

\begin{center}
  \uline{5:30 sáng\\Phòng của Kuon.}
\end{center}

Cứ tưởng rằng sau việc thức như ``cú đêm'' lúc một giờ sáng, tôi có thể ngủ một mạch đến
chín, mười giờ, thậm chí đến trưa, nhưng hỡi ôi, hóa ra đó chỉ là hy vọng hão huyền. Một
giọng nói trong trẻo, tinh nghịch vang lên bên tai, kèm theo những cái thúc nhẹ vào vai:
``Kuon-senpai, dậy đi, trời sáng rồi!''

Tôi lẩm bẩm trong cơn ngái ngủ, giọng pha chút bực dọc: ``{\color{red} Giời ơi, mấy giờ rồi mà lại om sòm thế\dots}''

Tôi lóng ngóng nghe được người đang phá hỏng giấc ngủ của tôi: ``\dots Ayame-chan?''

``Mặt trời sắp lên rồi đó! Dậy đi, Kuon-senpai!'' có vẻ như nàng Quỷ sừng đang bám lấy
vai, liên tục lay tôi dậy.

``\dots Mặt trời lên thì có liên quan gì tới anh đâu\dots\ Cho anh ngủ thêm chút nữa đi\dots'' mắt
tôi vẫn nhắm nghiền, nhưng tôi hoàn toàn có thể cảm nhận được sự bất mãn từ em ấy. Em
ấy phồng má, thở hắt ra đầy khó chịu.

Chắc không thể chịu nổi cái sự lười nhác của tôi, Ayame-chan đột nhiên làm một điều
khiến bố tôi, người ngủ cạnh, cũng phải giật mình tỉnh giấc.

Một tiếng \textbf{*KENG!}* chói tai vang lên. Tôi bật dậy hốt hoảng, thấy Ayame đang cầm một
cái dui gõ mạnh vào cồng chiêng - mà tôi còn chẳng biết nó được đặt ở đâu trong nhà.

``\textbf{Dậy đi, dậy đi, dậy!}'' em ấy vừa đập chiên vừa la lớn khắp cả căn phòng. Nghe thấy
tiếng chiêng cực đại, tôi giật mình, lắp bắp: ``Rồi rồi, anh dậy! Anh dậy đây!''

Ngay bên cạnh tôi, tiếng của bố tôi vọng ra, rõ ràng vẫn còn ngái ngủ: ``{\color{red} Có chuyện gì mà
  toáng cái mồm ghê thế!?}''

Ayame-chan nắm lấy tay tôi, cố gắng kéo tôi khỏi chiếc đệm: ``Dậy chuẩn bị ngắm mặt
trời mọc đầu năm (初日の出, hatsu hi no de) nào!''

Tôi nhíu mày, vẫn chưa hoàn toàn hiểu: ``Anh tưởng chỉ có Tết Dương Lịch ở Nhật Bản
mấy đứa mới có truyền thống đó thôi chứ? Người Etsunan bọn anh có làm chuyện này
đâu\dots''

Lập tức, Quỷ sừng-chan nhanh nhảu đáp lại với nụ cười rạng rỡ: ``Bình minh ngày đầu
năm ở Etsunan cũng đẹp lắm á! Nhanh lên nào! Yo!''

Tôi lê lết rời khỏi giường, cố gắng tập trung khi mắt còn lờ đờ. Thực ra, tôi đã thức đến hai
giờ sáng để\dots\ làm gì đó không quan trọng, nên bây giờ chẳng khác nào ``xác sống''. Cực
hiếm khi tôi có thể thức đêm đến thế, thậm chí thức qua đêm đến tận sáng.

Dưới ánh đèn mờ nhạt của mùa đông, tôi còn nghi ngờ trời hôm nay liệu có nắng nổi hay
không. Nhưng chưa kịp nghĩ xong, nàng Quỷ sừng đã lôi tôi từ phòng ngủ, qua hành lang,
xuống cầu thang nhà tôi, rồi lại leo lên tầng sáu của quán cà phê bên cạnh nhà.

Cứ mỗi lần em ấy bước xuống cầu thang, cả người tôi lại đập vào từng đó bậc gốm, rồi
lại va vào tường, thậm chí còn đụng cả lan can cầu thang, khiến cả thân người đau điếng.
Trời ạ, mới là buổi sáng sớm đầu năm thôi đấy\dots

``Ayame-chan à, tha cho anh đi mà\dots'' tôi kêu lên, trong khi đầu óc đang cố hiểu xem
chuyện gì đang xảy ra với bản thân.
Nhưng có vẻ như cô nàng Oni hớn hở đó không thèm đếm xỉa tới lời tôi nói: ``Nhanh lên,
Kuon-senpai ơi! Mặt trời sắp lên rồi!''

Dưới tầng, mẹ tôi - vẫn trong dáng vẻ bình thản thường ngày - nhìn hai chúng tôi, bật
cười: ``Ayame-chan dậy sớm ghê nhỉ. Cháu định ngắm mặt trời với mọi người à?''

Tôi quay lại, như vớ được chiếc phao cứu sinh: ``{\color{red} Mẹ ơi, cứu con\dots}'' Nhưng thay vì giúp tôi,
bà ấy chỉ mỉm cười, nhẹ nhàng giao phó: ``Thôi, cháu cứ lên đi. Shion-chan và Haato-chan
cứ lên với em ấy. Cô lên nhà chuẩn bị đồ ăn sáng. Nhớ trông chừng thằng Cò giúp
cô nha!''

Tôi trợn mắt, quay sang nhìn hai cô nàng đã trông nồi bánh chưng từ đêm qua. Hai người
lúng túng, đáp lí nhí: ``D-Dạ\dots''

Và thế là, với sự ``phản bội'' từ chính mẹ tôi, tôi đành bất lực nhìn nàng Quỷ sừng tiếp tục
kéo tôi đi, dù lòng không khỏi tự hỏi: ``Tại sao mọi chuyện lại thành ra thế này\dots''

Chúng tôi - mà thực ra chỉ có Ayame-chan, vì tôi bị em ấy lôi xềnh xệch như cái bao tải -
cuối cùng cũng chạy lên được tầng thượng của quán cà phê. Tất cả các thành viên \hll\
đã tề tựu đông đủ, trông ai cũng hào hứng chờ đợi.

``Bọn em đến rồi đây!'' em reo lên đầy tự hào, còn tôi thì ngã phịch xuống nền xi măng,
nằm bẹp dí, thở dốc. Nhưng niềm vui hội ngộ của mọi người chẳng kéo dài được lâu, vì
phần lớn ánh mắt nhanh chóng đổ dồn về phía tôi - người bị kéo lê đến nỗi trầy trụa, bầm
dập, và giờ thì còn không thể đứng dậy nổi.

Tôi thều thào nói với mọi người: ``Đừng lo cho anh\dots\ Chỉ là\dots\ cần thêm thời gian để
kiểm tra xem có gãy xương hay không thôi- Á!''

Một tiếng *RẮC* vang lên. Trời đất, chắc là tôi gãy mất cái xương sườn rồi. Cả đám
chạy đến vây quanh, vừa lo lắng hỏi han, vừa trừng mắt nhìn Ayame.

\begin{dialogue}
  \speak{Choco} Kuon-sama, anh ổn chứ?
  \speak{Kuon} \dots Ngoài việc anh không thể đứng dậy và có vài vết xô xát do Quỷ sừng kia kéo lê như cái chổi lau nhà, còn lại anh vẫn ổn.
  \speak{Aki} Ayame-chan, sao em phải kéo anh ấy như thế chứ?
  \speak{Rushia} Thật đấy! Anh ấy có giống như chúng ta đâu mà\dots
  \speak{Matsuri} Đúng rồi! Kuon-senpai đâu phải thánh thần hay quỷ dữ mà hành hạ như
  vậy!
  \speak{Subaru} Cô bị ngốc thật à!?
\end{dialogue}

Tôi bật cười khổ, nhắm mắt lại, chậm rãi đáp: ``Anh cứ cho đây là câu an ủi vậy\dots\ Mà
anh có chân mà, sao phải kéo lê anh thế\dots''

Trong khi đó, Akutan đã nhanh chóng bước đến, cúi người đỡ lấy tôi: ``Để em giúp anh
đứng dậy\dots'' Cùng với sự trợ giúp từ Choco-sensei, họ từ từ nâng tôi lên. Tôi phải dựa
vào hai người để không ngã lần nữa.

``Anh không nghĩ là bình minh này đáng giá đến thế đâu\dots'' tôi lẩm bẩm, làm mọi người
cũng chỉ cười trừ.

Tôi khẽ nhìn sang nàng Oni, em ấy cười gượng, tay xoa gáy, miệng lí nhí: ``Xin lỗi anh
nhé\dots\ Nhưng em chỉ muốn mọi người kịp xem bình minh thôi mà\dots''

Nhìn dáng vẻ lúng túng của Ayame-chan, tôi cũng không nỡ trách mắng thêm. Cuối cùng,
chỉ có thể lắc đầu cười trừ.

Thấy tôi đã chịu đứng dậy nhưng vẫn còn dựa dẫm vào hai người dìu, mọi người đều an
tâm phần nào. Họ cùng nhau bắt đầu hướng mắt lên trời, chờ đợi những tia mặt trời đầu
tiên trong ngày mùng Một Tết.

\dots Và đúng như dự đoán của tôi, trời vẫn u ám, mây đen dày đặc che kín, không một tia
sáng nào lọt qua nổi.

``Rõ ràng thế này mà vẫn bắt anh dậy\dots\ Mệt thật sự,'' tôi ngáp một hơi dài, lẩm bẩm ca
thán. Tôi biết rõ trời buổi sáng hôm nay sẽ âm u rồi, nên tôi không ngạc nhiên lắm.

Nhưng, một vài người trong các cô gái lại tỏ vẻ hoang mang hơn hẳn. Nhất là Matsurichan:
``Lạ nhỉ, giờ này rồi mà tại sao mặt trời vẫn chưa lên?''

A-san đứng cạnh liền giải thích với vẻ điềm tĩnh quen thuộc: ``Tôi nghĩ là do trời có quá
nhiều mây, Natsuiro-san.''

Mio-chan thở dài: ``Thế mọi người không ai xem dự báo thời tiết à?'' Okayu-chan và
Korone-chan gãi đầu cười trừ, tỏ vẻ ngượng ngùng:

\begin{dialogue}
  \speak{Okayu} Toàn tiếng bản địa, em nghe không hiểu\dots
  \speak{Korone} Em thì ngoài tiếng Nhật ra chẳng hiểu được ngôn ngữ nào hết á\~{}
\end{dialogue}

Sói đen-chan ngán ngẩm: ``Hai người có thể xem dự báo trên điện thoại được mà\dots''

Sora-chan ngước nhìn lên bầu trời, mặt đượm buồn thất vọng: ``Nếu trời cứ thế này thì
chắc ta không thể ngắm mặt trời mọc rồi.''

Tôi cố giữ chút tỉnh táo cuối cùng để nói trong cơn ngái ngủ: ``\dots Ai đó có thể vác anh về được không? Cho anh ngủ thêm tí nữa\dots''

Subaru-chan đập nhẹ vào vai tôi, tưởng rằng tôi sẽ lịm đi: ``Tỉnh dậy đi anh ơi!'' Fubu-chan
cũng nói góp: ``Chưa kể rằng sắp đến giờ mặt trời lên rồi đấy!''

Dù vậy, ngoài mấy người như A-san, Sora-chan, và Mio-chan tỏ ra khá bình thản, còn lại
thì giống như một đàn ong vỡ tổ, ai nấy đều lo lắng và cuống cuồng khi họ có thể bỏ lỡ
bình minh đầu năm.

``Cứ thế này thì chắc không ngắm được mặt trời mất\dots'' Aki-chan chạy tán loạn, răng bấm
móng tay em ấy liên tục. Tôi lập tức lên tiếng trấn an: `` Anh nghĩ mấy đứa không cần
tuyệt vọng vì chuyện này đâu. Chỉ là buổi bình minh thôi mà.''

Ayame-chan lập tức lên tiếng: ``Nhưng nếu không ngắm được, bọn em sẽ gặp xui đấy!''
Tôi thảng thốt ngạc nhiên, như thể vừa tỉnh dậy sau cơn sảng: ``Mê tín quá đó!!''

Trong lúc mọi người còn đang rối bời, đột nhiên, Pekora-chan reo lên: ``Sao chúng ta
không nhờ Miko-chan nhỉ, peko?''

Tất cả mọi ánh mắt lập tức đổ dồn về phía Vu nữ-chan, khiến chính em ấy cũng ngơ ngác:
``T-Tôi á, ne!?''

Peko-chan tiếp lời: ``Bà có cái quạt gì đó đúng không, peko?'' Roboco-chan cũng lập tức
hùa vào: ``Đúng đó, mình nhớ cậu từng dùng nó để thực hiện màn ``hoa vũ'' mà.''

``N-Nhưng cái đó thì giúp được gì, ne!?'' em ấy hoang mang, không hiểu họ đang ám chỉ
cô với điều gì.

``Chị có thể dùng chiếc quạt đó để thổi đám mây đi đó!'' Noel-chan nói, gợi ý cho hướng
đi của tiền bối mình.

``N-Nhưng mà quạt của tôi bé thế này, làm sao mà thổi chúng đi được chứ, ne! Đừng có
mà vớ vẩn!!'' Vu-nữ chan đáp lại, tỏ vẻ quả quyết rằng em sẽ không làm.

Nhưng rồi, dưới áp lực của cả nhóm, ai nấy đều đồng thanh hối thúc: ``Làm đi!'' liên tục,
Miko-chan cuối cùng cũng phải hét lên: ``\textbf{Thôi được rồi!}''

Em ấy bước ra mép tòa nhà, tay cầm chiếc quạt, rồi bắt đầu thực hiện màn nhảy gọi Thần
- hay ít nhất, mọi người vẫn gọi vậy. Tôi đứng đó, vừa ngạc nhiên vừa cố nhịn cười, vừa
thầm nghĩ: ``\textit{Mình không nghĩ em ấy sẽ làm được đâu. Ý tưởng hay đấy, nhưng chắc phải
  cỡ quạt Ba Tiêu\footnote{Món bảo bối của Bà La Sát (hay Thiết Phiến Công chúa), một nhân vật phản diện trong tác phẩm Tây Du Ký của Ngô Thừa
    Ân.} mới làm được.}''

Trong lúc nhảy, tôi còn nghe thấy Miko-chan còn lầm rầm như đang đọc thần chú: ``Kami-sama,
xin hãy nghe lời thỉnh cầu của bọn con!''

Kết thúc điệu nhảy, cả nhóm im lặng chờ đợi, ánh mắt đổ dồn lên bầu trời. Không có gì
thay đổi cả.

``Đã bảo rồi mà, thôi, đừng tốn công phí sức nữa. Mau đưa anh về đi-''
Nhưng rồi, như một phép màu, gió mạnh thổi qua, xua tan đám mây đen dày đặc, để lộ
một khoảng trời quang đãng ngay vị trí mà mặt trời chuẩn bị ló dạng.

``C-Cái quái!?'' tôi há hốc mồm trong sự ngỡ ngàng. Trong khi đó, các cô gái trố mắt nhìn
cảnh tượng kỳ diệu trước mặt, rồi cùng nhau vỗ tay rầm rộ, miệng không ngừng khen
ngợi.

Haato-chan reo lớn, ôm chầm lấy tiền bối của mình: ``Chị làm được rồi đó, Miko-chan!''
Trong khi đó, đến chính người đã thực hiện nghi thức đó vẫn còn đang bối rối: ``C-Chỉ là
hên (may mắn) thôi mà, ne!''

Tôi vẫn chưa tin được Miko-chan có thể làm được như thế.

Shion-chan thở dài nhẹ nhõm: ``May quá, đỡ cho tôi phải làm phép thuật rồi.'' Nghe vậy,
tôi nhếch mép: ``Lại lười chứ gì. Nhưng thôi, anh nghĩ là chúng ta đã nhờ em hơi nhiều
từ hôm qua rồi\dots''

Hiệp sĩ-chan tỏ vẻ phấn khích tột độ, khi khoảng trời trước mắt chúng tôi đã dường như
trở nên thoáng đãng: ``Như vậy là chúng ta có thể ngắm mặt trời mọc, đúng không!?''

Em ấy vừa dứt lời, Quỷ sừng-chan bỗng reo lên, chỉ tay về phía chân trời: ``\textbf{Mọi người
  ơi, nhìn kìa}!''

Mọi ánh mắt lập tức dõi theo hướng tay của em áy. Quả nhiên, từ phía xa, mặt trời đang
lấp ló nhô lên, ánh sáng đầu tiên của năm mới bắt đầu rọi xuống.

Ban đầu, chỉ là một viền sáng nhạt mỏng manh, nhẹ nhàng tô điểm cho những đám mây
gần chân trời. Nhưng chỉ trong vài phút, vầng dương dần nhô cao hơn, tràn đầy sức sống.
Màu sắc chuyển đổi ngoạn mục, từ sắc vàng kim óng ánh đến những tia cam rực rỡ, pha
thêm chút đỏ ửng của ngày mới.

Bầu trời phía trên cũng thay đổi, từ một màu xám u ám vì mây dày, giờ đây đã rực rỡ và
trong trẻo hơn bao giờ hết. Cảnh bình minh này, dù nằm giữa một thành phố đông đúc,
vẫn mang lại cảm giác thanh bình và thiêng liêng đến lạ thường.

``Đẹp ghê\dots'' Fubuki-chan khẽ thốt lên với khung cảnh bình minh hùng vĩ đó.

Ánh sáng bắt đầu lan tỏa, chiếu sáng những tòa nhà cao tầng xung quanh, làm nổi bật lên
từng góc cạnh như được dát vàng. Phía dưới, bóng cây và những con phố vắng lặng cũng
dần được chiếu sáng, một dấu hiệu rõ rệt rằng mùa xuân đã thực sự về.

Aki-chan mỉm cười nhẹ nhàng, tỏ vẻ mãn nguyện với những gì mà em, và chúng tôi,
được chiêm ngưỡng: ``Ayame-chan nói đúng\dots\ Bình minh ở đây cũng đẹp không kém gì
ở Tokyo cả.''

Thậm chí tôi, một người không trân trọng lắm những khoảnh khắc này, cũng phải gật đầu
đồng tình: ``Anh đồng ý.''

Những làn gió buổi sớm mang theo hơi lạnh cuối cùng của mùa đông dần tan biến, nhường
chỗ cho không khí trong lành và ấm áp hơn.

``Mình có thể cảm nhận được mùa xuân đang đến gần rồi đấy,'' Haato thốt lên, như thể
cảm xúc của em ấy đang dần hòa vào thiên nhiên vậy.

Nhưng, Ru-chan thì đã bắt đầu nhíu mày và lấy tay che mắt vì ánh nắng mặt trời lóe ra
quá sáng: ``Nhưng Rushia thấy hơi nóng rồi\dots\ Chói quá đi-nanodesu!''

Tiếng cười vang lên râm ran, xua tan những mệt mỏi từ đêm trước. Tôi bất giác đứng im,
để ánh sáng ấy chiếu lên gương mặt, cảm nhận hơi ấm từ mặt trời đầu năm mới.

A-san nhẹ giọng hỏi tôi, miệng chị ấy nở một nụ cười: ``Tôi nghĩ mọi người cũng đã thấy
thỏa mãn với cảnh này rồi. Hikawa-san, cậu cảm thấy sao?''

Tôi mỉm cười, đáp: ``Lâu lắm rồi em mới được chứng kiến cảnh này. Nó làm em nhớ lại
những ngày xưa\dots\ Nhưng chị đã quay lại cảnh này chưa ạ?''

Chị ấy đáp gọn nhẹ, làm tan biến đi những sự lo lắng của tôi với khung cảnh hùng vĩ này:
``Tất nhiên là ghi rồi.''

Chúng tôi cứ thế đứng đó, không ai nói gì thêm, chỉ chăm chú ngắm nhìn vầng dương
ngày càng nhô cao, ánh sáng lan tỏa khắp mọi nơi. Mặt trời cuối cùng cũng vượt qua ranh
giới của đường chân trời, trọn vẹn hiện ra, rực rỡ và hứa hẹn cho một năm mới đầy hy
vọng.

\ct

Sáu giờ ba mươi sáng ngày mùng Một Tết.

Bầu trời bên ngoài dần sáng rõ, những tia nắng đầu tiên của năm mới len lỏi qua khung
cửa sổ, tạo thành từng vệt sáng ấm áp trên sàn nhà gỗ.

\dots Tôi vẫn phải nhờ Subaru-chan và Aqua-chan dìu lên tầng ba. Những bước chân chậm
rãi của chúng tôi vang lên trên cầu thang gốm đá đã được chúng tôi lau dọn sạch sẽ, hoà
vào không gian tĩnh lặng buổi sáng sớm.

Phòng ăn đã sẵn sàng. Trên bàn là mâm cơm sáng với bánh chưng vuông vức, hành muối
chua ngọt, và một vài món ăn kèm truyền thống khác mà mẹ tôi đã chuẩn bị từ sớm. Mùi
thơm của bánh chưng, kết hợp với hương vị thanh mát của hành, khiến bụng tôi bất giác
kêu lên.

Bố mẹ tôi đã ngồi ở đó, dáng vẻ thư thái nhưng không giấu được nét tò mò khi thấy cả
đoàn đông đúc tiến vào.

Mẹ tôi mỉm cười dịu hiền: ``Chào buổi sáng mấy đứa.'' Bố của tôi cũng mỉm cười chào
mọi người, nhưng rồi tỏ vẻ ngạc nhiên khi nhìn tôi với đầy vết bầm tím trên nửa thân
dưới: ``Chân con bị làm sao thế?''

Tôi thở dài, giọng không giấu được chút hờn dỗi: ``Bố đi hỏi cái tên ``thủ phạm'' kéo con
xuống nhà ban nãy rồi lại kéo con lên nhà đi!''

Vừa dứt lời, ``thủ phạm'' mà tôi nhắc đến - nàng Oni tinh nghịch - thong thả bước vào
phòng, mặt tỉnh bơ như không có gì xảy ra. Ngay lập tức, các cô gái khác bắt đầu trách
mắng nàng vì đã làm tôi bầm dập'' hồi sáng. Đáp lại chỉ là cái nhún vai cùng một câu nói
đầy vẻ bông đùa: ``Loài người đúng là chân yếu tay mềm thật đấy.''

Fubu-chan cũng chỉ cười khổ: ``Không thay đổi gì cả\dots''

Trong khi đó, Subaru-chan, vẫn giữ vẻ mặt nghiêm túc, liền kể thêm với ông ấy về chuyện
sáng nay, còn Aqua-chan thì xoa bóp chân tôi, thi thoảng còn khẽ gật gù với những gì mà
bạn của em ấy nói.

\begin{dialogue}
  \speak{Subaru} Cậu ấy cũng thúc mọi người dậy hệt như thế, như thể đây là lần đầu tiên cậu ấy thấy bình minh vậy!
  \speak{Aqua} Cháu cũng không nghĩ Ayame-chan nhiệt huyết đến vậy đâu\dots
  \speak{Subaru} Thật đấy ạ.
\end{dialogue}

Bố tôi chỉ lắc đầu cười nhẹ, nhưng không quên quay sang quở trách tôi: ``Biết rõ như thế
rồi mà sao con không phản ứng gì?''

Tôi chỉ nhăn nhó, thở dài bất lực: ``Con còn đang ngái ngủ, thậm chí còn chả biết con là
đây, đâu là ai, rồi thế là Ayame-chan kéo con xuống ạ\dots''

Mẹ tôi nhìn cô nàng Quỷ sừng đang vui đùa với mọi người, khẽ cười trừ: ``Chả trách lại
bị thương như thế. Thôi, để tí nữa mẹ với Aqua-chan bôi thuốc cho.''

Tôi gật gầu, tỏ vẻ mệt mỏi: ``Vâng ạ\dots''

Dần dà, mọi người lần lượt ngồi vào chỗ. Không khí trở nên sôi nổi hơn khi những chiếc
bánh chưng được mở ra. Những tiếng cười vang lên khi từng chiếc bánh, vốn mang ``dấu
ấn'' riêng của từng người gói, được phát hiện.

Tôi cầm một cái bánh chưng trông có vẻ được gói tươm tất lên, nhưng rồi vết vá trên nó
khiến tôi nhận ra ngay: ``Cái bánh này hình như là của Akutan đúng không?''

``Ể!?'' chính chủ của chiếc bánh đó tròn mắt lên. Tôi giơ chiếc bánh, nói lớn: ``Thấy
có lá dong vá vào là biết mà.'' Tiếng cười rộ lên khi nàng Hầu gái đỏ mặt, vội núp sau
Shion-chan- để tránh ánh mắt trêu chọc.

Một âm vang bất ngờ khác lại thốt lên, khi lần này đến từ Aki-chan. Tôi nhìn thấy em
ấy đang cầm một chiếc bánh cỡ vừa, nhưng hai tay lại trùng xuống. ``Cái bánh này nặng
quá!''

``Lại còn có cả mảnh kim loại ở đây nữa,'' Mel-chan nhặt một sợi dây gai thép từ nó lên.
Quỷ sừng-chan nghe vậy, đắc ý đáp: ``À, cái đó là của tôi nha!''

Tôi ngán ngẩm nhìn em ấy: ``Vẫn nhất quyết gói kiểu `ném chó chó chết' à\dots''

A-san nhặt lên một chiếc bánh với những hình tròn và chữ X lớn, giơ lên hỏi: ``Bánh này
là của ai đây?'' Cặp đôi Chó-Mèo đồng loạt giơ tay: ``Của chúng em ạ!''

Mẹ tôi phì cười: ``Các cháu coi gói bánh chưng như một thú tiêu khiển đúng nghĩa đen
nhỉ.'' Lời trêu chọc khiến cả bàn cười vang, chỉ trừ Pekora-chan, người vừa thở dài vừa
lắc đầu đầy ngán ngẩm: ``Tôi đã nói thế rồi mà hai người không nghe, peko\dots''

Roboco, với dáng vẻ ngượng ngùng, vội vàng nhận lỗi khi chiếc bánh của em ấy ``to đến
mức bà T*n vê-lốc vái em ấy bằng cụ'': ``Xin lỗi mọi người nha, tại em hăng quá\dots'' Chiếc
bánh phải đến hơn hai mươi người mới ăn hết.

Tôi thở dài, giọng tỏ vẻ châm chọc: ``Ít nhất còn có một PON biết gói bánh chưng đúng
cách đấy.'' Nghe vậy, Roboco đưa tay che mặt đang đỏ hỏn lại: ``Đừng khịa em nữa đi\~{}!''

Bữa sáng diễn ra trong không khí sôi nổi. Mọi người vừa thưởng thức bánh chưng, vừa
bàn tán về kế hoạch cho ngày hôm nay. Ánh nắng buổi sớm chiếu rọi qua cửa sổ, làm
sáng bừng cả căn phòng, mang theo hy vọng cho một năm mới đầy khởi sắc.

\ct

\pov{-}

Buổi sáng mùng Một Tết, sau bữa ăn sôi nổi cùng đại gia đình tại tầng ba nhà Hikawa,
nhóm chúc Tết chuẩn bị khởi hành. Theo lời đề xuất của Kuon, nhóm gồm hai bố con nhà
Hikawa - Ryujin và Kuon - cùng bốn thành viên \hll : nàng Ma cà rồng Mel, nàng
Thần tượng Sora, và bộ đôi Cáo trắng - Sói đen Fubuki và Mio.

Với các cô gái, họ mặc bộ trang phục truyền thống mà họ mang tới, được thể hiện và có
mô tả từ hình dưới đây:

\begin{figure}[H]
  \centering
  \begin{subfigure}[t]{0.2\textwidth}
    \centering
    \includegraphics[height=8cm]{../../../../Image/Book?/newyear/sora_model.png}
    \caption{Sora diện một bộ
      furisode (có thể nhận
      diện bằng tay áo dài và
      được làm bằng lụa rất
      mịn, màu sắc tươi sáng
      hơn) màu xanh, tay rộng
      với hoạ tiết hoa trên đó.
      Cô có một dây buộc obi
      (dây buộc quanh lưng)
      màu vàng với một dây
      thừng buộc màu đỏ
      (obiage) cùng với nó. Cô
      có cài một chiếc trâm
      bông hoa đỏ, gọi là
      kanzashi.}
  \end{subfigure}
  \quad
  \begin{subfigure}[t]{0.2\textwidth}
    \centering
    \includegraphics[height=8cm]{../../../../Image/Book?/newyear/mel_model.png}
    \caption{Mel diện một bộ
      kimono hai màu: đỏ ở
      trên và tím ở dưới, có
      họa tiết là các hình ngôi
      sao, váy hakama dài đến
      cổ chân màu tím, đi tất
      tabi trắng và dép geta.
      Bằng cách nào đó, cô có
      kiểu tóc dài hơn, được
      buộc thành hai bím tóc
      thấp với phần tóc mái
      được trang trí bằng nơ đỏ
      và kẹp tóc.}
  \end{subfigure}
  \quad
  \begin{subfigure}[t]{0.2\textwidth}
    \centering
    \includegraphics[height=8cm]{../../../../Image/Book?/newyear/fubuki_model.png}
    \caption{Fubuki diện một bộ
      kimono trắng, chuyển
      dần sang tông màu xanh
      nhạt khi xuống dưới, một
      chiếc váy hakama cạp
      cao màu xanh, một chiếc
      obi màu xanh và đi bốt
      buộc dây màu đen. Tóc
      của cô được trang trí
      bằng một chiếc trâm
      kanzashi làm từ bông
      hoa đỏ, và có một nút
      thắt hoa có tua rua.}
  \end{subfigure}
  \quad
  \begin{subfigure}[t]{0.2\textwidth}
    \centering
    \includegraphics[height=8cm]{../../../../Image/Book?/newyear/mio_model.png}
    \caption{Mio diện một bộ
      kimono màu đỏ, trên cổ
      đeo một chiếc vòng đen,
      có tay áo rộng in hoa,
      váy hakama ngắn đến đùi
      màu đen, quần tất đen và
      đi bốt buộc dây màu nâu.
      Khi ở nhà, cô có mũ đội
      đầu của hầu gái và có tạp
      dề eo trắng xếp ly, đồn
      thời tay áo cũng ngắn
      hơn. Cô để tóc của mình
      thành đuôi ngựa, lệch
      hẳn sang bên trái.}
  \end{subfigure}
  \label{setone}
\end{figure}

Với hai bố con nhà Hikawa, họ mặc bộ trang phục tươm tất và lịch sự hơn so với ngày
thường: Kuon mặc áo sơ mi sọc xanh có cúc bấm, mặc một chiếc quần bò xanh có thắt
lưng màu nâu và đi giày tây đen, còn Ryujin thì diện một chiếc áo len màu đen cùng với một
chiếc quần bò đen có thắt lưng màu đen và đi giày thể thao màu xanh sẫm.

Thông thường, nếu chỉ có hai bố con, hai người họ sẽ đèo nhau bằng xe máy và đi đến
các nhà để chúc Tết. Nhưng năm nay, gia đình cậu làn này chơi trội hơn: họ mượn một
chiếc ô tô của mẹ cậu, chở cả sáu người họ đi, do bố cậu (cực hiếm khi) cầm lái. Chiếc xe
rời khỏi khu phố lúc gần tám giờ sáng, mọi người vừa đi, vừa trò chuyện với nhau đầy hồ
hởi. Kuon ngồi ghế trước, vẫn giữ vẻ mặt hơi lơ mơ vì còn ngái ngủ và ê ẩm thân người.
Phía sau, bốn cô gái vừa ngắm khung cảnh bên ngoài vừa không ngừng bàn tán về không
khí Tết.

Mel háo hức: ``Chị không ngờ ở Etsunan lại có không khí Tết truyền thống đậm chất
như vậy!''

Mio mỉm cười: ``Chỗ này đúng là yên bình thật. Không giống như những nơi chúng ta
hay làm việc\dots\ Chắc giờ này ở Tokyo đông nghịt lắm đây.''

``Thì bởi vì ở Nhật Bản hôm nay vẫn là ngày bình thường mà,'' Kuon lên tiếng, ``họ bây
giờ có nghỉ Tết Nguyên Đán như bọn anh đâu.''

Fubuki gật gù đồng ý với hai người: ``Phải ha. Nhưng mà này, em không thấy chỗ nào
bán đồ ăn vặt trên đường cả. Sao em thấy thiếu thiếu nhỉ?''

Sora, với một người đã nghiên cứu kỹ lưỡng về ngày Tết trong vài ngày qua, nhẹ nhàng
giải thích: ``Vì mùng Một người ta không bán hàng đâu.''

Kuon khẽ bật cười, tiếp lời: ``Thế nên anh mới bảo là phải đi mua đồ từ trước là vậy đó.
Bây giờ gần như chẳng có hàng quán nào mở cửa giờ này đâu.''

Nghe các cô gái và cậu con trai của ông bàn luận: ``Thế thì mấy đứa cứ thoải mái trải
nghiệm đi. Nhà nội của Kuon đón Tết cũng `chất' lắm đấy.''

\ct

Ngôi nhà của một người họ hàng nhà nội Kuon nằm sâu trong một con ngõ nhỏ, bao quanh
bởi những ngôi nhà mặt tiền to lớn. Khi chiếc xe vừa đỗ lại, cả nhóm nhanh chóng xuống
xe, chỉnh trang trang phục trước khi bước vào bên trong ngõ. Bên trong con ngõ hơi tối tăm nhưng đã được dọn dẹp sạch sẽ đó là đầy rẫy những ngôi
nhà khác. Nhà của một người họ hàng có thể phân biệt được bằng hai lớp cửa: một chiếc
cửa kéo tay bằng nhôm ở bên ngoài, và một cánh cửa kính đẩy ở bên trong.

Cả hai bố con mở cánh cửa kéo ra, rồi cùng nhau đẩy chiếc cửa kính vào, lộ ra hai người
đang ngồi trên ghế bành trắng đang chờ đợi bọn họ: một người đàn bà lớn tuổi, tóc đã
điêm những chấm bạc, và một người đàn ông trông trẻ hơn Ryujin, đầu gần như trọc lốc.
Người đàn ông đó cũng là người làm việc lâu năm với ông đã được gần mười năm.

``{\color{red} Bà, bọn cháu đến chúc Tết đây ạ,}'' người bố nói, theo sau là cậu Tổng chào hai người
bằng phép lịch sự tối thiểu: ``{\color{red} Cháu chào bà, cháu chào chú Min(h).}''

Các cô gái cũng đồng loạt cúi đầu chào, khiến cả hai người họ bất ngờ. Người phụ nữ lớn
tuổi nhìn hai bố con cậu, rồi quay sang các cô gái, ánh mắt lộ vẻ tò mò: ``{\color{red} À, đây chẳng
  phải là bạn của thằng Cò sao?}''

Ryujin gật đầu, giải thích: ``{\color{red} Dạ, đây là Mel-chan, Sora-chan, Mio-chan và Fubuki-chan
  - đồng nghiệp làm việc chung với Kuon. Các cháu lần đầu đến Etsunan nên con mời về
  chơi dịp Tết.}''

Người đàn ông nhếch mép cười, nhìn cậu Tổng rồi giở giọng trêu: ``{\color{red} Thằng Kuon này cũng
  gớm phết đấy! Được làm việc với gái như thế sướng quá còn gì!}''

Câu đùa của người đàn ông khiến cậu con trai đỏ mặt, còn các cô gái thì cười rúc rích.

\ct

Tại mỗi nhà mà nhóm ghé qua, không khí đều náo nhiệt và vui vẻ. Những câu chuyện
trao đổi xoay quanh năm mới, công việc, và đặc biệt là sự xuất hiện của bốn cô gái từ
\hll.

Họ nhanh chóng chiếm được cảm tình của mọi người, đặc biệt là Mel, với nụ cười duyên
dáng và cách nói chuyện nhẹ nhàng. Mio và Fubuki thì không ngừng bày trò đùa nghịch
với đám trẻ con, trong khi Sora lại ghi điểm bởi vẻ điềm đạm và chững chạc.

Sau mỗi lượt chúc Tết, cả nhóm lại nhận được những phong bao lì xì đỏ thắm. Fubuki,
không giấu được sự tò mò, liên tục mở từng phong bao để kiểm tra bên trong.

Đến khi kết thúc lượt chúc Tết ở nhà anh trai của Hijaki - tức anh trai trưởng của gia đình
Kuon, cô nàng than vãn: ``Có từng này thôi à\dots''

Mio nhíu mày nhìn cô bạn: ``Lì xì không phải là khoản thu nhập đâu, Fubuki. Còn chưa
kể bà phải mừng tuổi lại bọn trẻ con đấy.''

Nàng Dơi tỏ vẻ đồng cảm với cô, rồi nói với vẻ hãnh diện: ``Giống như chúng tôi làm hồi
sáng nay chẳng hạn.''

Sora mỉm cười: ``Đúng rồi. Quan trọng là cái lộc từ người ta trao cho mình đó, Fubukichan
à.''

Dù vậy, nàng Cáo trắng vẫn ca thán: ``Nhưng lộc như thế này thì ít quá! Làm sao mà đủ
để gacha đây\dots''

Kuon chỉ biết lắc đầu, bật cười nhẹ: ``Em cứ như bọn trẻ con bây giờ ấy. Thôi, sang nhà
tiếp theo đi!''

Họ tiếp tục ghé qua ít nhất năm nhà nữa trước khi kết thúc chuyến đi vào lúc gần trưa. Cả
nhóm quay về nhà Hikawa trong tâm trạng hào hứng, sẵn sàng cho bữa trưa đặc biệt với
món bánh mochi tự tay giã.

Buổi sáng đầu năm trôi qua êm ả, để lại trong lòng mỗi người những kỷ niệm đáng nhớ
về không khí Tết truyền thống và sự ấm áp từ gia đình Kuon.

\ct

\begin{figure}[H]
  \centering
  \begin{subfigure}[t]{0.2\textwidth}
    \centering
    \includegraphics[height=8cm]{../../../../Image/Book?/newyear/rushia_model.png}
    \caption{Rushia mặc bộ
      kimono màu đen họa tiết
      hoa anh đào, buộc dây
      obi màu vàng giữa lưng.
      Cô còn có tà áo trắng
      mang họa tiết con bướm
      tím. Cô búi tóc lệch trái,
      đeo nơ bướm lớn trên
      đỉnh đầu, đeo khuyên
      hình bướm ở tai trái, có
      tóc xanh chuyển dần tím
      hồng, đi tất tabi trắng và
      dép zouri đen quai đỏ.}
  \end{subfigure}
  \quad
  \begin{subfigure}[t]{0.2\textwidth}
    \centering
    \includegraphics[height=8cm]{../../../../Image/Book?/newyear/ayame_model.png}
    \caption{Ayame diện một bộ
      furisode màu đỏ có họa
      tiết hoa, buộc dây obi
      màu vàng, một chiếc
      haori (áo choàng) màu
      tím, đi tất tabi trắng và đi
      dép okobo màu đen. Cô
      để tóc rất dài được buộc
      thành bím tóc kiểu Pháp
      ở một bên và buộc đuôi
      ngựa ở bên còn lại, được
      trang trí bằng hoa
      kanzashi có tua rua.}
  \end{subfigure}
  \quad
  \begin{subfigure}[t]{0.2\textwidth}
    \centering
    \includegraphics[height=8cm]{../../../../Image/Book?/newyear/pekora_model.png}
    \caption{Pekora mặc một bộ
      kimono váy ngắn màu đỏ
      viền ren, tay áo xắn lên,
      Don-chan được dùng
      như một tasuki (ruy-băng
      buộc tay áo) để giữ, dây
      buộc obi màu vàng có nơ
      lớn ở sau lưng, tất tabi
      hình con thỏ màu trắng
      và dép okobo màu đen.
      Cô để búi tóc tết đôi
      được buộc bằng khăn
      kanzashi có tua rua.}
  \end{subfigure}
  \quad
  \begin{subfigure}[t]{0.2\textwidth}
    \centering
    \includegraphics[height=8cm]{../../../../Image/Book?/newyear/miko_model.png}
    \caption{Miko mặc kimono
      không tay xếp nếp, cổ áo
      rời hồng với ruy băng tua
      rua, tay áo rời so le trắng
      và hồng nhạt, váy cạp
      cao màu hồng, xếp nếp
      với nhiều lớp lót trong,
      đeo túi xách hình bánh
      cá, đi tất tabi trắng và dép
      geta đỏ. Tóc cô buộc bím
      kiểu Pháp và tết hai bên
      với băng đô màu trắng
      đính tua rua và kẹp tóc.}
  \end{subfigure}
  \quad
  \begin{subfigure}[t]{0.2\textwidth}
    \centering
    \includegraphics[height=8cm]{../../../../Image/Book?/newyear/noel_model.png}
    \caption{Noel mặc một bộ
      furisode có họa tiết hoa
      xanh lam xen lẫn
      yagasuri, thắt dây chiếc
      obi xanh lam với một
      cuộn obiage xanh lá cây,
      đi tất tabi trắng và đi dép
      geta.Cô để tóc đuôi ngựa
      ngắn với phần mái được
      chia đôi và trang trí bằng
      trâm cài kanzashi có gắn
      hoa màu trắng và xanh.}
  \end{subfigure}
  \quad
  \begin{subfigure}[t]{0.2\textwidth}
    \centering
    \includegraphics[height=8cm]{../../../../Image/Book?/newyear/roboco_model.png}
    \caption{Roboco diện một bộ
      kimono ngắn màu đỏ có
      tay áo xếp nếp, thắt dây
      obi màu đen có phần tay
      áo rời kéo dài đếc cổ tay
      có chung màu và hoa văn
      với kimono trang trí với
      một vài dây vàng. Cô để
      tóc đuôi ngựa ngắn buộc
      bằng nơ lớn màu đỏ và
      kanzashi từ bông hoa.}
  \end{subfigure}
  \quad
  \begin{subfigure}[t]{0.2\textwidth}
    \centering
    \includegraphics[height=8cm]{../../../../Image/Book?/newyear/okayu_model.png}
    \caption{Cô mặc bộ kosode
      trắng với tay áo rộng
      viền ruy băng màu tím,
      một chiếc váy dài bằng
      hakama tím xếp nếp, một
      chiếc tạp dề trắng in hình
      onigiri góc dưới bên phải
      cố định bằng một dải như
      obi đen, đi tất tabi trắng
      và dép okobo đen. Cô để
      tóc như mọi khi, cài trâm
      màu tím và đeo chuông
      tóc nhỏ.}
  \end{subfigure}
  \quad
  \begin{subfigure}[t]{0.2\textwidth}
    \centering
    \includegraphics[height=8cm]{../../../../Image/Book?/newyear/korone_model.png}
    \caption{Cô mặc kimono
      không tay hai tông màu
      trắng-xanh với dây buộc
      chéo, một chiếc khăn
      choàng hagoromo, tay áo
      rộng rời, đi tất tabi màu
      trắng và geta màu xanh.
      Cô để tóc rất dài với
      những chiếc kẹp tóc hình
      xương và tóc tết hai bên
      có ống tóc thông thường.}
  \end{subfigure}
  \label{settwo}
\end{figure}

Cùng thời điểm mà sáu người họ đi chúc Tết và nhận tiền lì xì, trên vỉa hè trước nhà của
cậu Tổng, không khí mùng Một cũng tràn đầy năng lượng chẳng kém.

Nhóm ở lại gồm Kaien, Ryuuhou, và một vài thành viên \hll\ như Rushia, Ayame,
Noel, Roboco, Pekora, Miko, Okayu và Korone. Công việc chính của họ là giã bánh
mochi từ những phần gạo nếp đã được ngâm kỹ từ đêm trước và hấp chín từ sáng sớm.

Người mẹ của gia đình đứng ở trung tâm, vừa chỉ đạo vừa chuẩn bị các công đoạn. Con
gái của nhà nhanh nhẹn giúp mẹ mình sắp xếp dụng cụ, từ chày, cối giã gạo, đến khay
đựng mochi đã hoàn thành.

Để hưởng ứng ngày Tết truyền thống, các cô gái cũng đã diện bộ trang phục truyền thống.
Với hai mẹ con nhà Hikawa, cô em gái Ryuuhou diện một chiếc áo dài có tà đỏ và đi một
đôi bốt trắng cao cổ, người mẹ Kaien thì mặc một chiếc áo dài màu hồng, khoác bên ngoài
một chiếc áo da để tránh rét. Còn về các thành viên của \hll\, cũng mặc bộ trang phục
truyền thống mà họ mang tới, cho thấy sự giao thoa giữa Tết truyền thống của Etsunan và Nhật Bản\footnote{Xem lại trang trước đó để xem.}.

Công việc tưởng chừng đơn giản, nhưng với sự góp mặt của những thành viên nổi tiếng
hay làm quá như Haato, Rushia, và Ayame, mọi chuyện nhanh chóng rẽ hướng sang\dots\ hỗn loạn đầy hài hước.

Người mẹ nhà Hikawa đứng giữa sân, chỉ tay về phía cối giã gạo bằng gỗ đã sẵn sàng:
``Mọi người chú ý nhé! Gạo nếp đã hấp chín rồi, ai khỏe thì giã, ai khéo thì vò bánh. Làm
nghiêm túc thì bánh mới ngon được!''

Ayame, với dáng vẻ nghịch ngợm thường thấy, liền giơ tay: ``Em xin xung phong! Em sẽ
là người quản lý gạo nếp!''

Cô em gái của Kuon thắc mắc: ``Ayame-onee-san, chị chỉ cần xúc gạo thôi mà cũng tranh
giành ạ?''

``Vậy thì cô em cũng làm phụ Rushia-chan đi chứ nhỉ, yo?'' nàng Quỷ sừng nói. Ryuuhou
thở dài, đi về chiếc cối mà Rushia đã chờ sẵn.

Hai người họ sẵn sàng cầm chày, dù cây chày gỗ trông gần như quá nặng so với cơ thể
nhỏ bé của họ. Cô em gái nhướng mày nhìn cộng tác của cô: ``Rushia-pnee-chan, hôm nay
phải giã mạnh vào nhé! Onii-san không muốn nhìn thấy chị yếu đuối đâu nhỉ?''

Rushia đỏ mặt, vừa ngượng vừa nổi cáu với đàn chị của mình: ``C-Cái gì chứ!? Chị không
có yếu! Đưa chày đây, chị sẽ giã nát hết mọi thứ cho mà xem!''

Cặp đôi bắt đầu công đoạn giã bánh. Ban đầu, mọi thứ có vẻ chậm rãi và đều nhịp\dots

\begin{dialogue}
  \speak{Ryuuhou} Rushia-onee-chan, đập mạnh nữa vào! Mạnh nữa!
  \speak{Rushia} C-Chị không nghĩ là cái chày này lại nặng thế\dots
  \speak{Ryuuhou} Ơ? Sao em tưởng chị còn hùng hổ lắm cơ mà?
\end{dialogue}

\dots nhưng rồi ngay khi Ayame đi ngang qua hai người họ, buông một câu đầy tính châm
chọc: ``Rushia-chan, nghĩ về việc \textit{GAMERS} đang tranh giành Kuon-senpai đi!'' Lập tức,
cô trố mắt ra, sực nhớ lại vụ việc ngày hôm qua. Ngay sau đó, cô biến thành một ``cô đồng
yandere'' chính hiệu.

``Rồi, lại nào!'' cô em gái của Kuon vừa vò xong bánh mochi, không để ý mọi thứ xung
quanh. Tức khắc, Rushia đã đập xuống cối với tốc độ cực nhanh, trong khi không ngừng
la hét: ``\textbf{Lũ khốn kiếp!!}''

Với một tiếng hét vang trời, cô nâng cây chày gỗ nặng như bấc rồi giáng xuống cối gạo,
khiến cả sân vang lên tiếng *bốp* mạnh đến nỗi Ayame phải giật mình làm rơi cả vá xúc
gạo, còn Ryuuhou phát hoảng lùi ra phía sau.

Pekora trố mắt ngạc nhiên: ``R-Rushia!? Bà định giã bánh hay đang phá cái cối vậy?''
Nhưng dường như nàng Pháp sư không để tâm tới mọi người, mà tiếp tục vừa đập chày
vào cối vừa la hét.

``B-Bình tĩnh lại nào, Rushia-onee-san! Chị đập như thế kẻo vỡ cối đó!'' Ryuuhou ôm lấy
nàng tóc xanh đang điên cuồng kia, nhưng dường như không hiệu quả, mà thậm chí còn
làm cho cô nàng ấy mạnh hơn: ``\textbf{Đây là số phận của kẻ dám động vào Kuon-senpai!!!}''

Trong lúc đó, nàng Quỷ sừng thì cười không ngớt, vừa liên tục hò hét cổ vũ Rushia: ``Đúng
rồi, mạnh nữa đi! Tưởng tượng như em đang đối đầu với Korone-chan ấy!''

Korone, đang nhai dở miếng bánh mẫu bên cạnh, lập tức thốt ra ngạc nhiên: ``Ể? Sao lại
lôi tôi vào đây vậy? Măm măm\~{}''

Okayu, người cũng đang ăn thử miếng mochi khác, lên tiếng động viên: ``Thôi nào Korosan.
Ít ra bà còn an toàn hơn cái cối gạo kìa, măm măm\~{}''

Trông nàng Pháp sư đang giã bánh một cách điên loạn, Noel từ phía xa nhìn cô một cách
ngưỡng mộ: ``Mình không nghĩ là mình có thể giã bánh nhiệt huyết tới như vậy\dots''

``Đừng có mà chọc tức người ta nữa, peko!'' nàng Thỏ trong khi đó thốt lên tỏ thái độ bất
bình, trong lúc tay còn đang mải vò bột bánh với Noel.

Đến cả cặp đôi Roboco và Miko cũng ngừng tay khi nhìn cặp Ryuuhou và Rushia đang
giã bánh.

\begin{dialogue}
  \speak{Noel} Rushia-chan lúc ghen trông mạnh thật đấy\dots
  \speak{Miko} Thật là kinh khủng, ne\dots
\end{dialogue}

Nàng Người máy vốn đã khá lúng túng khi phải giã gạo nặng tay, không may nhìn chằm chằm cặp
Ryuuhou và Rushia đang làm loạn bên cạnh. Cô cầm cây chày một cách lơ đễnh, không
để ý Miko đang cúi xuống sửa lại miếng vải bọc cối.

``Roboco, cẩn thận tay chày của bà kìa, không là đập trúng tay chị bây giờ đó, ne!'' Đối
phương cười khì, tỏ vẻ rất tự tin: ``Ể? Đừng lo, Miko-chan, tôi kiểm soát được mà!''

Nhưng vừa dứt lời, nàng Người máy vung mạnh cây chày, và\dots\ *bốp!* Cây chày đập
trúng vào tay phải của nàng Vu nữ. Cô hét toáng lên: ``\textbf{Áaa! Tay tôi, ne\~{}!!!}''

Cả sân im lặng vài giây, rồi Kaien hoảng hốt chạy tới kiểm tra. ``M-Miko-chan! Cháu có
sao không vậy!?'' May mắn thay, dù cho nàng Người máy đập chày khá mạnh, cô chỉ bị
đau nhẹ. Dù thế, cô chắp tay trước ngực, giả vờ than thở, mắt chảy lệ: ``Chết rồi, giờ sao
tôi ăn bánh được đây, ne\~{}?''

Nàng Thỏ thì không tiếc lời mà mắng tiền bối của mình: ``Đồ hậu đậu, Roboco-senpai!
Làm việc gì cũng peko-éo xong nổi, peko!!''

Roboco gãi đầu cười gượng, trong khi Kaien đặt băng vải lên tay Miko rồi nhìn sang cô,
nửa đùa nửa thật: ``Có khi cô sẽ cho cháu cầm cái chày nhựa thay vì cái chày gỗ đấy,
Roboco-chan ạ.''

Ở một góc khác, Rushia dường như mất kiểm soát hoàn toàn, mắt long lên sòng sọc. Cô
không chỉ giã mạnh xuống cối của mình, mà còn chuyển mục tiêu sang\dots\ cối của cặp
Okayu-Korone ở ngay bên cạnh, khiến hai bọn người hoảng loạn.

``N-Này, Rushia-chan, em làm cái gì vậy!?'' nàng Mèo bám lấy một cánh tay của nàng
Pháp sư. ``Đừng có phá hỏng bột bánh của bọn chị chứ!'' nàng Khuyển bám lấy cánh tay
còn lại.

``Xin lỗi hai chị ạ, tại Rushia-onee-chan bỗng nhiên phát điên\dots'' Ryuuhou vẫn đang
cố bám thân trước của nàng.

Cả ba người, dù đã cố hết sức để ghìm chặt Rushia, cô vẫn thẳng tay đập xuống vừa nhanh,
lại vừa mạnh. Tiếng bột nếp văng tứ tung làm hai cô nàng rên rỉ:

\begin{dialogue}
  \speak{Okayu} Tạm biệt, bánh mochi của mình\~{}\dots
  \speak{Korone} Rushia-chan! Bọn mình có làm gì cậu đâu mà lại tàn nhẫn với bọn này thế!
\end{dialogue}

Rushia tiếp tục điên cuồng lấy cây chày đập thẳng xuống nền vỉa hè, lao thẳng tới cối
giã của cặp đôi Pekora-Noel. Nàng thỏ thấy người bạn đồng thế hệ của mình đang điên
cuồng, liền chạy lại cố gắng ngăn cản.

Nhưng khi vừa nắm tay Rushia, nàng Thỏ bị Pháp sư đẩy ra một cách mạnh mẽ, trong khi
Ryuuhou thì bị đẩy xuống ngã phía sau.

\begin{dialogue}
  \speak{Pekora} N-Này! Tôi chỉ định giúp bà thôi mà, peko!
  \speak{Ryuuhou} Rushia-onee-chan, dừng lại đi! Đừng làm hỏng mọi thứ nữa!
\end{dialogue}

Cuối cùng, kẻ đầu têu cho mọi chuyện - Ayame - đành phải lấy bánh mochi tươi ra dụ dỗ:
``Rushia-chan, ăn đi! Nhìn này, Kuon-senpai cũng thích bánh mochi thế này lắm đấy!''

Tới lúc này, mặt của nàng Pháp sư mới giãn ra, tỏ vẻ vui tươi, sáng lấp lánh như sao: ``Ơ!?
Thật sao? Đưa em ngay!!''

Cô ngắm nhìn bánh mochi một cách rất kỹ lưỡng, dường như đã hạ nhiệt. Tất cả mọi
người đều thở phào nhẹ nhõm, riêng nàng Quỷ sừng vẫn nở một nụ cười tinh nghịch.

\ct

Trong lúc các nhóm khác đang làm việc, nàng Oni, với tính cách nghịch ngợm vốn có,
tiếp tục nghĩ ra một ý tưởng điên rồ khác. Cô chạy tới chỗ Noel, thì thầm điều gì đó vào
tai nàng Hiệp sĩ. Nàng tóc xám bật cười: ``Ý kiến hay đấy, làm luôn đi!''
Nàng Quỷ sừng liền trèo lên lưng nàng Hiệp sĩ, cầm chày giã bánh từ trên cao giáng xuống
cối gạo, tạo nên cảnh tượng như ``chiến binh khổng lồ''.

\begin{dialogue}
  \speak{Noel} Lên đi nào, Ayame-senpai!
  \speak{Ayame} Xem bọn này đây, kỹ thuật giã bánh mạnh mẽ nhất thế giới! \textbf{Gyaaah!}
\end{dialogue}

\dots Tất nhiên, ``cỗ máy'' Ayame-Noel hoạt động không hiệu quả chút nào. Chày vừa giã
xuống, bột mochi bắn lên tung tóe, dính đầy mặt cả hai. Kaien đứng gần đó chỉ biết thở
dài: ``Ayame-chan, Noel-chan, hai đứa có định làm việc nghiêm túc không vậy?''

Cả hai cô nàng phủi bột khỏi mặt, rồi cười trừ trước mặt người mẹ của gia đình.

``Cháu xin lỗi cô ạ, Kaien-san\dots'' nàng Hiệp sĩ cười khì, trong khi nàng Quỷ sừng thì cười
gian: ``Bọn cháu đang thử nghiệm phương pháp mới thôi ạ!''

Miko, người vẫn đang ngồi chườm đá gần đó, thảng thốt: ``Phương pháp giã bánh mới cái
con khỉ, ne!''

\ct

Sau khi được băng bó tay, nàng Vu nữ vẫn không chịu nghỉ ngơi. Cô nằng nặc đòi giã
bánh tiếp, nhưng lần này để bản thân cô tự giã, thay cho Roboco.

``C-Chị có chắc là chị cầm được cái chày gỗ không vậy, Miko-onee-chan? Trông nó còn
to hơn cả chị đấy\dots'' Ryuuhou lên tiếng tỏ vẻ lo lắng cho cô. Nàng Vu nữ cười một cách
khoái chí: ``Có gì đâu, cũng giống như chị giã bánh mochi đợt trước thôi mà, ne!''

Hai mẹ con nhà Hikawa thấy vậy, cũng chỉ thở dài tạm thời an tâm, khi ít nhất nàng Vu
nữ đã có kinh nghiệm giã bánh mochi từ trước.

\dots Để rồi sự an tâm liền vụt tắt ngay khi cô cầm chiếc chày. Ngay khi Miko-chi nâng nó
lên bằng cả sự can đảm, cô vung chày giã thật mạnh\dots\ nhưng chẳng may lại đập trúng
vào chiếc xô đựng nước đặt gần đó.

``Ối!!! Nước đâu mà tràn ra thế này, ne\~{}!?'' nàng Vu nữ hoảng loạn. Dòng nước chảy
lênh láng làm cả sân trở thành một chiến trường trơn trượt. Korone mất thăng bằng, ngã
ngửa xuống đất, lôi theo Okayu. Pekora cũng không tránh khỏi, trượt chân và ngồi bệt
xuống.

Trong khi đó, Ayame thì lại cười khoái chí với khung cảnh hỗn loạn này, không quên tranh
thủ quay video: ``Cảnh này đáng giá lắm đây! Đúng là đầu năm đầy may mắn!''

Kaien thì chỉ để tay đặt lên trán, nhưng vẫn mỉm cười: ``Cô nghĩ là sau Tết này nhà cô phải
mua thêm bảo hiểm nhà mới được\dots''

\ct

Trong khi mọi người vẫn đang cố gắng dọn dẹp bột nước tung tóe khắp sân, Ayame nhảy
lên một bậc thềm gần đó, hô to: ``Thông báo! Cuộc thi giã bánh mochi chính thức bắt
đầu! Ai giã nhanh và đẹp nhất sẽ được thưởng\dots\ một phong lì xì từ Kaien-san!''

Mẹ đang bận lau sạch đôi dép bị dính nước, giật mình ngẩng lên: ``Ơ từ từ, cô chưa chuẩn
bị tiền lì xì-'' nhưng chưa kịp dứt lời, nàng Quỷ sừng tinh nghịch đã bồi thêm một câu đầy
sát thương: ``Và phần thưởng đặc biệt\dots\ là một đêm ngủ với Kuon-senpai!''

\il

Trong khi đó, Kuon - ``giải thưởng đầy bất đắc dĩ'' này - đang trên đường trở về nhà sau
chuyến đi chúc Tết. Bất chợt, cậu hắt hơi một tiếng: ``\textbf{Hắt xì!}''

Cả bốn cô gái lẫn bố cậu tỏ vẻ ngạc nhiên. ``Sao thế, Kuon(-senpai)?''

``Tự dưng con cảm thấy lạnh quá\dots'' cả người cậu khẽ run lên cầm cập, rồi lại hắt hơi lần
nữa.

``Lạ nhỉ, bố có bật điều hòa đâu nhỉ?'' Ryujin chỉnh vặn nút tìm công tắc điều hòa, trong
khi bốn cô gái thì tỏ vẻ ngạc nhiên.

``Kuon-senpai, anh ốm sao?'' Mel lên tiếng hỏi, tỏ vẻ lo lắng. Cậu vội đáp: ``K-Không
phải thế\dots''

``Kể cũng lạ, hôm nay trời có lạnh lắm đâu nhỉ?'' Mio tỏ vẻ băn khoăn, Sora cũng gật đầu:
``Ừ nhỉ, hôm nay dự báo thời tiết nói hôm nay trời đẹp thích hợp để ra ngoài mà?''

``Không lẽ nào\dots'' Fubuki nói, giọng của cô nàng tỏ rõ sự gian xảo, ``có ai đó đang nhắc
tới anh sao?''

Bạn của cô vỗ nhẹ vào lưng, nhíu mày: ``Đừng có mà vớ vẩn nữa, Fubuki! Tôi không
nghĩ anh ấy sẽ hắt hơi vì cái chuyện ba xàm này đâu-''

Bất chợt, Kuon lên tiếng chen ngang: ``Không\dots\ Anh nghĩ là Fubu-chan nói đúng đấy.''
Cả bốn cô nàng nhìn cậu, vừa hoang mang, vừa nghiêng đầu khó hiểu.

``Ý anh là sao, Kuon-senpai?'' Sora hỏi. Cậu đáp: ``Anh đang có cảm dự không lành ở
nhà\dots''

Mặc cho bố cậu nói như trách mắng: ``Vớ vẩn! Con cứ nghĩ linh tinh!'' nhưng Kuon vẫn
không giấu nổi sự lo lắng. ``Không biết ở nhà đang diễn ra chuyện gì nhỉ?''

\il

Cả sân như đông cứng lại trong vài giây. Các cô gái quay đầu nhìn Ayame, rồi bỗng dưng
mặt của họ đỏ hỏn hơn hẳn.

``M-Một đêm ngủ chung với Kuon-senpai sao?'' tất cả bọn họ đồng thanh. Để rồi\dots

\dots Rushia, như một quả bom nổ chậm, bùng nổ ngay tức khắc. Cô hét lên đầy phấn khích:
``Cái gì!? Một đêm với Kuon-senpai ư!? Em sẽ thắng cuộc thi này!'' Noel cũng chen vào,
vỗ tay hưởng ứng: ``Ý hay lắm, Ayame-senpai! Để xem ai có thể giành được giải thưởng
này nào!''

Pekora đứng bên cạnh, không khỏi lườm nàng Oni: ``Mới sáng ra đã bày trò rồi, peko!''
rồi lại tỏ vẻ phấn khích: ``Nhưng\dots\ một đêm với Kuon-senpai, nghe cũng không tệ!''

Okayu và Korone nhìn chung có vẻ khá bàng quan với giải thưởng này. ``Bọn mình đã
ngủ chung với anh ấy rồi, đúng không Ogayu\~{}'' nàng Chó hỏi. Cô bạn trả lời: ``Đúng vậy
á\~{} Chắc lần này mình sẽ chỉ tham gia cho vui thôi\~{}''

Riêng chỉ có Kaien ôm trán, thở dài nhưng vẫn nở một nụ cười: ``Đúng là có Ayame-chan
thì nhà cô có một ngày Tết độc nhất vô nhị thật!''

\ct

Roboco vẫn còn hơi lúng túng sau tai nạn với Miko, nên cô quyết định sử dụng chế độ
``siêu tốc'' để giành chiến thắng. Cô bật chế độ tăng lực, khiến cánh tay máy của mình
xoay như chong chóng.

``Mọi người, nhìn đây! Công nghệ vượt trội sẽ chiến thắng!''

Tuy nhiên, không biết có phải do bản năng ``Đại ngốc'' của cô vẫn còn ở đó, hay do cô
chưa hoàn thiện chế độ này, mà chày gỗ trên tay cô đã văng thẳng vào cối của cặp đôi
Okayu-Korone lần nữa, làm bột bắn lên tung tóe.

Nàng Miêu thấy đống bột của hai người lại rơi vãi xuống, khẽ cười trừ: ``Roboco-senpai
thật sự không muốn bọn này làm xong bánh trong bình yên thật à\~{}'' Cô bạn Khuyển thì
bật cười: ``Em nghĩ chị nên dừng lại trước khi cả sân này thành chiến trường á, Roboco-senpai!''

Trong khi đó, Rushia, với quyết tâm cao độ, tiếp tục giã bánh mạnh cùng với Ryuuhou
tiếp tục khéo tay vò bánh.

``Đập chậm thôi chị ơi, em còn vò nữa chứ!'' cô em gái kêu lên, nhưng đáp lại cô là một
giọng hối thúc như một con bò tót Tây Ban Nha: ``Phải làm thật nhanh, làm thật ngon, thì
ta mới dành chiến thắng!''

Chiếc chày giã vào rất mạnh, và rất nhanh, tới mức mà cối của họ gàn như không chịu
nổi. Ryuuhou thở dài: ``Nếu chị cứ đập mạnh thế này, coi chừng chị làm hỏng luôn cái
cối đấy!''

Nàng Pháp sư lườm cô em gái khiến cho đối phương phát sợ, rồi bất ngờ\dots\ đưa cho cô
một cái chày gỗ khác. ``Em có giỏi thì thử thi với chị một phen!''

Ryuuhou lớ ngớ cầm lấy cái chày gỗ, rồi nhìn nàng Pháp sư đang vò bánh một cách ngạc
nhiên. ``Còn đứng ngây ra làm gì! Làm đi!''

``V-Vâng!''

Dường như vì cảm nhận được sự đáng sợ từ Cô đồng đó mà cô em gái của Kuon đã liên
tục đập mạnh vào chiếc cối giã gạo, trong khi Rushia thì tiếp tục vò bánh với vẻ mặt quyết
tâm.

Các nhóm khác nhìn hai người họ giã bánh và vò bánh đều tỏ ra ngỡ ngàng, thậm chí còn
đổ đồ mồ hôi hột. Trong khi đó, Ayame, đứng xem từ xa, vừa cười vừa trêu: ``Hai cô bé
này có vẻ tập trung nhầm đối tượng rồi thì phải!''

Ở một góc khác, Pekora nảy ra ý định tinh nghịch. Cô lấy một chiếc bánh mochi nhỏ còn
sống, nặn thành quả cầu tuyết và nhắm thẳng vào Ayame. ``Xem đây, peko! Đây là phần
thưởng cho trò đùa của chị!!''

Nhưng không may, nàng Quỷ sừng nhanh nhẹn tránh được, không quên cất lên: ``Ố, hụt
rồi!'' Quả cầu bánh mochi đó bay thẳng vào mặt của Miko. Cô la lên oai oái: ``P-Pekora!
Em đang làm cái gì vậy ne!?''

Cả sân lại cười nghiêng ngả, trong khi Pekora lén lút quay lưng chạy trốn: ``Ơ, xin lỗi, xin
lỗi! Không cố ý đâu, peko!''

Trong lúc cô đang loay hoay tìm đường ``trốn án'', Noel vô tình vung cây chày lên để
chuyển sang cối khác. Nhưng vì hơi mất tập trung, chày gỗ của cô vô tình đập trúng ngay
đầu của nàng Thỏ đang hốt hoảng kia.

*BỐP!*

Chiếc chày tác động vào đầu quá mạnh khiến nàng Thỏ lăn quay một vòng rồi ngã vật
xuống đất. Nàng Hiệp sĩ vội bỏ chiếc chày lại, rồi đi tới cô nàng vừa bị đập kia: ``Á! Pekora-chan, mình xin lỗi! Cậu có sao không!?''

Nàng Thỏ xanh chỉ kịp kêu lên một tiếng yếu ớt: ``Peko\dots'' và ngất lịm đi, khiến mọi
người hoảng loạn chạy tới. Ayame vội nắm lấy tay Kaien, vừa cười vừa cố nhịn: ``Không
sao đâu, chắc em ấy chỉ ngất xỉu tí thôi. Nhưng mà\dots\ Noel-chan, mạnh tay thật đấy!''

Nhìn mọi người đang hớt hải dìu nàng Thỏ vào trong nhà, Kaien chỉ có thể đứng cười trừ,
đùa: ``Có cần cô mua bảo hiểm cho từng đứa không đây?''

Cuối cùng, đích thân bà phải làm thay cho Pekora, đội sẵn mũ bảo hiểm để an toàn.

\ct

Sân nhà giờ đây trông như một bãi chiến trường thực sự. Những mảnh vụn gạo vương
vãi khắp nơi, cối gạo nghiêng ngả, vài chiếc chày nằm chỏng chơ bên cạnh bên trong nhà
để rửa.

Mặc dù vậy, không khí tràn ngập tiếng cười và sự hài lòng. Cuối cùng, mọi người cũng
hoàn thành công việc của mình. Những khối bánh mochi mềm mại, trắng mịn được xếp
ngay ngắn trên những khay lớn, còn tỏa hơi ấm của sự chăm chút và tình bạn.

Cũng đúng lúc này, đội chúc Tết cũng đã trở về nhà, mong chờ món mochi do những
người trong đội làm. Không chỉ có hai bố con nhà Hikawa, mà còn bốn cô gái cũng rất
mong chờ.

Kaien bước lên bậc thềm nhỏ bên cạnh cây đầo rực rỡ. Trong bộ áo dài truyền thống màu
vàng tươi, bà nở nụ cười dịu dàng, đôi tay nhẹ nhàng phủi lớp bột dính trên vạt áo trước
khi cất lời: ``Cảm ơn mọi người đã tham gia nhiệt tình! Dù có chút\dots\ hỗn loạn.''

Bà nhìn thoáng qua chiếc cối bị Noel làm gãy trước đó, rồi tiếp tục nói: ``cô vẫn rất vui
vì tinh thần đồng đội của cả nhóm. Nhìn những khay bánh đẹp đẽ kia, cô chắc chắn đây
sẽ là một kỷ niệm đáng nhớ. Và bây giờ, hãy thưởng thức thành quả của chúng ta nào!''

Với việc bà là một MC chương trình hội chợ ca nhạc, nên giọng nói của bà gây rung động
được trái tim cho mọi người. Mọi người đồng loạt vỗ tay, tiếng cười nói râm ran khắp sân.
Kuon, với bộ kimono giản dị (cậu bất đắc dĩ phải thay do chính cậu cũng không biết vì
sao), ngồi xuống chiếu trải giữa sân cùng các thành viên. Trước mặt cậu là bốn đĩa bánh
mochi được trình bày công phu bởi các đội. Mỗi đĩa đều có một phong cách riêng, phản
ánh rõ cá tính và ``sức sáng tạo'' của từng đội trong buổi sáng hỗn loạn vừa qua.

Cậu cầm đôi đũa lên, nhìn bốn đĩa bánh, nhíu mày nhìn thành quả của cả bốn đội. ``Vậy
thì\dots\ ta bắt đầu từ đâu nhỉ?''

Cậu vừa nói vừa nhìn quanh, đôi mắt hiện rõ vẻ vui thích pha lẫn tò mò. Các cô gái lập
tức xôn xao:

\begin{dialogue}
  \speak{Pekora} \hl{Từ đội bọn tôi chứ, peko! Đội chúng tôi chắc chắn là đẹp nhất mà, peko!}
  \speak{Okayu} \hl{Đừng tự tin quá, ha\~{} Bánh của tụi chị ngon hơn, chỉ cần Kuon thử là biết ngay!}
  \speak{Rushia} \hl{Kuon-senpai, bắt đầu với đội em đi! Đó là\dots{} toàn bộ trái tim và linh hồn của em đấy!}
\end{dialogue}

Mẹ của Kuon bật cười khẽ, nhẹ nhàng ngồi xuống cạnh cậu con trai: ``Thôi nào, từ từ nào
các cô gái. Kuon ăn thử từng món một, ai cũng sẽ có cơ hội.''

Câu nói của bà làm dịu bớt không khí cạnh tranh, dù ánh mắt của vài người vẫn lộ rõ sự
quyết tâm. Kuon nhấc đũa lên, bắt đầu hành trình thưởng thức từng chiếc bánh mochi,
sẵn sàng đối mặt với những điều bất ngờ từ mỗi đội.

Ngay khi cậu đang chuẩn bị thưởng thức bánh mochi, Ayame đột ngột đứng lên, tay giơ
cao một chiếc chày nhỏ, giọng đầy phấn khích: ``Khoan đã! Còn một điều nữa mà Kuon-senpai chưa biết!''

Tất cả ánh mắt đổ dồn về nàng Quỷ sừng, trong đó có cả sự tò mò, lẫn lo lắng, dù rằng họ đã được đoán được ý của cô nàng.

``Thật ra thì\dots\ bọn em đã vừa tổ chức làm bánh mochi đó, Kuon-senpai! Mình đã tổ chức
nó với một giải thưởng đặc biệt!''
Kuon nheo mắt, bắt đầu cảm thấy có điều gì đó không ổn: ``Giải thưởng? Em đang nói
cái gì vậy?''

``Đừng lo, Kuon-senpai, phần thưởng rất hấp dẫn!'' nàng Quỷ sừng nháy mắt với cậu, rồi
tiếp tục hướng về phía mọi người: ``Ai làm bánh ngon nhất sẽ được\dots''

Cả sân im lặng chờ đợi câu trả lời, Ayame dừng lại một chút để tăng phần kịch tính, rồi
tung ra một câu nói khiến tất cả đứng hình: ``\textbf{\dots một đêm ngủ chung với Kuon-senpai!}''

Ngay lập tức, Kuon giật mình đến mức suýt ngã ngửa ra sau. Cậu lắp bắp: ``C-Cái gì cơ!?
Ayame-chan em đang đùa đúng không!?''

A-chan đứng gần đó, chỉ lắc đầu thở dài: ``Nakiri-san, tôi tưởng Tết là thời điểm để chúng
ta làm những điều tốt đẹp, chứ có phải làm thêm náo loạn đâu\dots''

Trong khi đó, các cô gái, dù đã biết được giải đặc biệt đó, vẫn không giấu nổi sự phấn
khích, với mỗi người một cách khác nhau:

\begin{dialogue}
  \speak{Rushia} Tuyệt vời! Giải thưởng này là của em chắc luôn!
  \speak{Okayu} Ohoho\~{} Nghe cũng thú vị đấy, nhưng liệu em có thể vượt qua được đội tụi này không?
  \speak{Pekora} Không phải ngủ chung với ai đó là phần thưởng kỳ cục nhất mà tôi từng
  nghe sao, peko!? Cơ mà\dots\ nó cũng hay đấy!
  \speak{Korone} Thế mà cũng nghĩ ra được!
  \speak{Miko} Đã bảo là Ayame-chan sẽ lại bày trò gì kỳ quặc mà, ne!
\end{dialogue}

Dưới ánh mắt đầy hoang mang của Kuon, Ayame chỉ cười khúc khích, không hề tỏ ra hối
lỗi. Trong khi đó, Kuon thở dài, chấp nhận thực tế rằng một ngày yên bình sẽ không bao
giờ tồn tại khi ở cạnh các cô gái này.

\ct

Pekora vẫn còn ôm đầu đau nhức sau cú chày ``trời giáng'' của Noel, đành ngồi một góc
thở dài bất lực. Kaien nhẹ nhàng thay cô đảm nhiệm vị trí, đôi tay khéo léo tạo hình những
khối mochi trắng ngần. Khi bánh được dọn lên, Kuon lấy đũa gắp một miếng, đưa lên
miệng cắn thử.

Cậu nhai một cách chậm rãi, thưởng thức vị ngọt của chiếc bánh. Cậu khẽ gật gù: ``Mẹ,
bánh của mẹ với Noel-senpai làm rất ngon. Vị ngọt vừa đủ, dẻo mềm đúng chuẩn luôn
ạ.''

Bất chợt, cậu nhướng mày, quay sang mẹ mình hỏi: ``Nhưng mà\dots\ mẹ có cần phải làm tới
mức đánh ngất đồng đội không ạ?''

Cả nhóm cười rộ lên, khiến Noel đỏ mặt, vội vàng xua tay phân bua: ``K-Không phải đâu!
Đ-Đây chỉ là tai nạn thôi mà!'' Cô lúng túng cúi đầu xin lỗi người đồng đội của mình:
``Pekora-chan, mình xin lỗi nhé! Mình không cố ý đâu mà, thật đấy!''

Pekora khoanh tay, mặt nhăn nhó, nhưng vẫn không giấu được vẻ hài hước trong ánh mắt:
``Không cố ý mà lại chính xác đến thế à, peko!? Chắc chắn năm sau tôi sẽ không làm với
Noel nữa, peko!''

Tiếng cười lại vang lên khắp sân, làm xoa dịu đi phần nào những ``tổn thương'' của nàng
Thỏ. Mẹ của Kuon vừa cười vừa lắc đầu, thêm vào: ``Cô cũng nghĩ năm sau nếu các cháu
có tới thì nên chia Noel ra làm một nhóm riêng cho an toàn.''

Lúc này, Sora, Mel, Mio và Fubuki cũng không nén được ý kiến của mình. Nàng Thần
tượng lên tiếng nụ cười nhẹ, luôn giữ vẻ điềm tĩnh vốn có: ``Nhưng phải công nhận là đội
này làm bánh rất khéo. Nhìn khối mochi này, ai mà nghĩ được cảnh tượng trước đó nhỉ?
Với cả, chị thấy cách phối màu và độ mềm đều rất đồng đều đấy!''

Mel nghiêng đầu, tay cầm một miếng bánh, tỏ vẻ thích thú: ``Ừm, bánh này không tệ,
nhưng có lẽ vị hơi nhạt hơn chút xíu so với đội khác. Chắc tại em lo xử lý `đồng đội' hơn
là chăm chút công thức, đúng không, Noel-chan?''

Noel cúi mặt, lúng túng đáp: ``Em đã cố hết sức mà, mọi người đừng làm em xấu hổ nữa
chứ!''

Trong khi đó, Fubuki thưởng thức một miếng, nhắm mắt như đang suy ngẫm: ``Thật ra
chị thấy sự kiện `trận chiến' này rất phù hợp với không khí ngày Tết. Mấy đứa đúng là
hơi hỗn loạn nhưng cũng đầy niềm vui á\~{}''

Mio nhăn nhó, đập nhẹ vào lưng của bạn mình: ``Đó không phải là điều mà chúng ta cần
vào ngày Tết đâu, bà biết đấy\dots''

Cáo trắng cười tươi rói, nháy mắt về phía hai thành viên của đội thứ nhất: ``Cúng ta cũng được thêm vài kỷ niệm đáng nhớ nữa. Năm sau nhớ rủ cả Noel
và Pekora tiếp tục tham gia ha\~{}''

Pekora nhăn mặt, nhưng trước sự ủng hộ vui vẻ của mọi người, cô chỉ thở dài: ``Được rồi,
peko. Nhưng lần sau tôi phải được chọn đối tác an toàn hơn đấy!''

Cả sân lại bật cười sảng khoái, tiếng cười hòa quyện cùng không khí ấm cúng, nhộn nhịp
của buổi sum họp đầu năm. Cuối cùng, cậu Tổng chỉ biết cười trừ trước màn đối đáp đầy
sinh động. Cậu cầm thêm một miếng bánh, giơ lên: ``Thôi nào, mọi người. Dù có chuyện
gì thì đây vẫn là thành quả rất tuyệt vời. Bánh ngon, cả nhóm cố gắng hết mình. Như thế
là đủ rồi!''

Lời nói của Kuon khiến bầu không khí lắng dịu lại, và cả nhóm cùng mỉm cười, tận hưởng
những khoảnh khắc ấm áp bên nhau.

\ct

Tới lượt đội của Rushia và cô em gái của Kuon. Nàng Pháp sư ngồi cạnh Ryuuhou, ánh
mắt sáng rực đầy tự hào khi thấy Kuon cầm chiếc bánh mochi từ khay của đội mình. Lần
này cậu dùng tay trần để cầm lên ăn, thay vì đũa do cậu cho là không cần thiết.

Ngay khi Kuon chuẩn bị cầm lên một miếng, nàng Pháp sư đột nhiên nắm chặt tay cậu,
ánh mắt cháy bỏng: ``Kuon-senpai, anh phải ăn thử miếng đầu tiên! Đây là bánh mochi
mà em đã dồn hết tâm huyết vào đấy!''

Cậu ngơ ngác, nhìn bàn tay đang nắm lấy mình rồi quay sang Ryuuhou để định hỏi cho
ra nhẽ. Cô em gái chỉ nhún vai, vẻ bất lực: ``Em cũng không hiểu tại sao chị ấy lại nhiệt
tình thế, Onii-san. Chỉ biết là lúc thi đấu, chị ấy nói liên tục về `ý nghĩa của tình yêu đích thực' lắm ạ.''

Cậu con trai gãi đầu, vẫn chưa hoàn toàn hiểu: ``Ừm\dots\ nhưng đây chỉ là bánh mochi thôi
mà, có gì đặc biệt đâu?''

Câu nói khiến Rushia như bị đánh trúng tim, cô lập tức đứng bật dậy, giơ tay lên cao như
diễn thuyết: ``Không phải chỉ là bánh mochi! Đây là biểu tượng của nỗ lực, của đam mê,
của\dots\ của tình yêu! Anh không hiểu gì hết, Kuon-senpai à!''

Đằng sau Kuon, Sora khẽ bật cười, thì thầm với Mel: ``hlEm ấy nói cứ như thể đang thi
đấu `Vua đầu bếp' đấy nhỉ.'' Nàng Dơi cười khì, đáp: ``\hl{Mà, hình như Kuon-senpai vẫn
  chưa nhận ra, nhỉ? Anh ấy ngây thơ quá!}''

Quay lại với Cô đồng vẫn đang mải miết giảng giải cho món ăn của mình. Cậu cố gắng
xoa dịu bằng cách nhấc một miếng bánh lên: ``Rồi, rồi. Anh biết rồi. Đừng có kích động
quá. Để ăn thử xem nào.''

Vừa thấy cậu đưa bánh lên miệng, Rushia lập tức nghiêng người về phía trước, ánh mắt
đầy mong đợi: ``Sao rồi, Kuon-senpai? Ngon chứ!? Em đã làm hết sức vì anh đấy!''

Cậu nhai chậm rãi, gật đầu nhẹ: ``Ừm, vị dẻo mềm, ngọt nhẹ, rất vừa miệng. Phải nói là
rất ngon. Nhưng mà\dots''

Vừa lúc cậu nói ``nhưng'', nàng Pháp sư đột ngột nín thở, vẻ căng thẳng như đang chờ đợi
một lời nhận xét sống còn. Rushia và Kuon ngước về hai chiếc cối gạo bị đập hỏng ở góc
vỉa hè, nhướng mày hỏi: ``\dots em có nhất thiết phải đập hỏng cả hai cái cối gạo chỉ để làm
được bánh này không?''

Câu nói khiến cả sân bật cười. Ryuuhou gật đầu, đồng tình: ``Em phải công nhận là
Rushia-onee-chan hơi\dots\ quá sức. Nhìn chị ấy thi đấu mà em tưởng chị là một chiến binh
huyền thoại vậy, tới mức em giã còn không kịp tay mà.''

Rushia đỏ mặt, cắn môi, nhưng vẫn cố gắng biện minh: ``E-Em chỉ muốn mang lại món
bánh ngon nhất cho Kuon-senpai thôi mà!'' Cô hắng giọng, rồi nhanh chóng lấy lại phong
thái tự tin: ``Tất cả là vì anh, Kuon-senpai! Nếu cần, lần sau em sẽ làm lại với ba cái cối
cũng được!''

Cậu Tổng cười gượng, đưa tay lên trán như không biết phải phản ứng thế nào. Fubuki
nhìn cảnh tượng ấy, trêu chọc: ``Kuon-senpai, anh đúng là người có `sức ảnh hưởng' mạnh
mẽ thật đấy!''

Mio thì khẽ lắc đầu, nhìn cậu với ánh mắt đồng cảm: ``Đúng hơn là cậu ấy không có quyền
từ chối đâu. Nhìn ánh mắt bùng cháy của Rushia-chan đi là đủ hiểu.''

Mel gật đầu, cười trừ: ``Công nhận, nhưng mà cái giá phải trả cho món bánh này hơi\dots\ cao quá.''

Nàng Pháp sư quay đầu lại, khoanh tay và hất mặt đầy kiêu hãnh: ``Hãy gọi đó là `sự
hy sinh vì tình yêu' đi, Kuon-senpai!'' Tất cả mọi người nhìn nhau, không biết phải nói sao cho cô nàng
đang hăng máu kia. Ryuuhou chỉ biết lắc đầu ngao ngán: ``Em bắt đầu sợ chị rồi đó,
Rushia-onee-san\dots''

\ct

``Được rồi, đội tiếp theo.'' Okayu và Korone đẩy đĩa bánh của mình về phía Kuon. Những
miếng mochi trắng mịn, căng tròn, trông hoàn hảo như được làm bởi những đầu bếp
chuyên nghiệp.

Korone cười hớn hở, chỉ tay về đĩa bánh: ``Anh phải ăn thử ngay, Kuon-senpai! Đây là
công sức của tụi mình sau biết bao lần thất bại á!''

Cậu con trai nhìn cặp đôi với vẻ tò mò, rồi lấy một miếng bánh lên và nếm thử. Vừa cắn
một miếng, cậu đã gật gù, vẻ mặt đầy hài lòng: ``Hai đứa làm cũng tốt đấy. Vị rất mềm
mại, ăn cũng khá mịn. Rất chuẩn vị.''

Nàng Miêu nháy mắt tự hào, giọng ngọt như một cô mèo con: ``Em biết mà. Bọn em lúc
nào cũng là một đội hoàn hảo cả na\~{}'' Nàng Khuyển trong lúc miệng vẫn đang nhai bánh,
cười lớn: ``Tôi nghĩ năm sau chúng ta nên gắn biển cảnh báo `Nguy hiểm' quanh khu vực
làm bánh mất!''

Cậu Tổng nhíu mày, tỏ vẻ ngạc nhiên và khẽ bật cười: ``Sao lại phải cảnh báo? Hai đứa
đang nói gì vậy?''

Nàng Mèo bật cười, vỗ vai người bạn: ``Bọn em làm tốt sau khi cản phá hai lần ạ. Lần
đầu là do Rushia-chan đánh trúng cối bột của tụi chị trong lúc giã bánh.'' Nàng Chó minh
họa cho cậu bằng một đường chày quét ngang: ``Cối bột bay tung tóe luôn! tôi còn tưởng
đó là chiêu thức của một võ sĩ chuyên nghiệp nữa chứ!''

Okayu gật đầu, tiếp tục: ``Lần thứ hai thì là Roboco-senpai. Chị ấy nói `Chỉ thử một chút thôi,' nhưng lại tăng lực quá mạnh, làm cối thứ hai cũng đổ luôn.'' Korone phụ họa xoa
cằm, tỏ vẻ nghiêm trọng như đang phân tích một trận đấu: ``Đúng vậy! Roboco-senpai có
khi còn mạnh hơn cả Rushia-chan. Chỉ một cú là bột của bọn em đi tong lần nữa á!''

Cậu Tổng nhìn cặp đôi đang than phiền với biểu cảm vui vẻ đó, rồi khẽ nhìn nàng Pháp
sư và nàng Người máy - những người bị cặp đôi này réo tên - đang chỉ cười nhạt trước tội
của mình.

Trong khi đó, Sora bật cười, quay sang Mio và Mel: ``Không ngờ Okayu và Korone lại
gặp toàn `đối thủ' mạnh thế nhỉ?'' Mio thở dài, rồi cười mỉm: ``Đúng vậy. Nhưng dù sao
họ vẫn hoàn thành bánh, thế mới đáng khen.''

Mel mỉm cười, ánh mắt trêu chọc: ``Có khi nào đây là kế hoạch của Rushia-chan không?
Em ấy có vẻ như muốn cản đường tất cả để đội mình chiến thắng ha.''

Fubuki khẽ bật cười: ``Nếu thế chắc tôi cũng thừa sức làm được. Nhưng với tư cách là
trưởng nhóm \textit{GAMERS}, kết quả của hai người họ thật sự rất ấn tượng!''

Kuon mỉm cười, nhìn lại đĩa bánh: ``Như vậy là đĩa bánh của hai đứa được mọi người
đánh giá cao đấy. Nhưng mà, lần sau tránh xa tên đại PON đằng kia nha!''

Nàng Người máy nghe vậy, liền phồng má giận dỗi: ``Em không có PON mà! Mà tại sao
anh chỉ nhắm mỗi em không vậy\~{}?''

Câu nói của Kuon khiến cả nhóm phá lên cười. Korone và Okayu cũng không nhịn được
cười, cùng nhau cười nghiêng ngả:

\begin{dialogue}
  \speak{Korone} Chắc chắn lần sau bọn em sẽ chọn chỗ an toàn hơn!
  \speak{Okayu} Nhưng mà cũng vui mà, nhỉ\~{}?
\end{dialogue}

\ct

Chỉ còn lại đĩa bánh cuối cùng, đến từ cặp Miko và Roboco. Cậu Tổng nhìn đĩa bánh của
họ với chút tò mò. Những chiếc bánh không hoàn toàn tròn trịa mà có hình dáng hơi méo
mó, trông giống như những viên đá cuội kỳ lạ.

``Sao bánh của hai đứa trông méo vậy?'' cậu cất giọng hỏi trong sự ngờ vực. Lập tức,
nàng Vu nữ không ngần ngại chỉ vào đồng đội: ``Tại cậu ta đó, ne!''

``Ể? Sao lại là tôi chứ?'' Roboco trố mắt ra, ``Miko-chan cũng là người vò bánh mà?''

Dù vậy, cậu vẫn lấy một miếng lên và cắn thử. Cậu nhai chậm rãi, rồi bật cười nhẹ:
``Không tệ, nhưng hình dáng thì hơi\dots\ độc đáo nhỉ.''

Miko đứng cạnh Roboco, tay vẫn đang băng bó, mặt đỏ bừng vì ngượng lập tức phân bua:
``Tại Roboco làm rơi chày vào tay tôi, ne\~{}! Đừng có trách một mình em chứ!''

Nàng Người máy thì cúi đầu cười ngượng ngùng: ``Tôi không có nhớ là chế độ siêu tốc
này vẫn chưa hoàn thiện\dots\ Tôi đâu có biết đâu chứ!''

Nàng Vu nữ lắc đầu thở dài, giọng đầy vẻ phàn nàn: ``Bà đúng là hậu đậu thế là cùng, ne!
Chắc chắn năm sau tôi sẽ không làm với Roboco nữa đâu, ne\~{}! Lần nào cũng gặp chuyện kỳ lạ!''


``Nhưng bà cũng là một đại Ngố nữa trong công ty mà, Miko-chan,'' Roboco nói, xóc xỉa
với đồng đội mình. Lập tức, đối phương tức giận, trừng mắt về phía cô: ``Bà đang nói cái
gì đó, ne\~{}!? Bà mới là người hậu đậu ở đây đấy!''

Cả nhóm bật cười trước màn đối đáp của hai người. Sora nháy mắt với Mio, vừa cười vừa
nói: ``Tôi nghĩ đội này chắc là thú vị nhất đấy.''

Nàng Sói đen tỏ vẻ nghiêm túc hơn: ``Thú vị thì đúng, nhưng mà cũng nguy hiểm không
kém. Chỉ cần nhìn bánh của họ thôi là đủ biết rồi.''

Nàng Cáo trắng chỉ tay vào một chiếc bánh có hình thù kỳ quái, trông giống như một
chiếc mũ nhỏ: ``Bánh mochi này trông giống\dots\ chiếc mũ samurai! Có khi nào họ định
làm một bộ sưu tập không?''

Mel mỉm cười, thêm vào: ``Ừ, chắc đây là `phiên bản giới hạn' chỉ có đội Miko-chan và
Roboco mới làm được à.''

Kuon bật cười, đồng tình với hai người họ: ``Dù sao thì cũng rất sáng tạo. Vị bánh vẫn
ngon, dù hình dáng có hơi\dots\ khác biệt một chút.''

Nàng Vu nữ nghe vậy, tỏ vẻ đắc ý, khoanh tay tự hào, dù mặt cô vẫn còn hơi đỏ hỏn: ``Đó!
Thấy chưa! Dù sao cũng là bánh ngon mà!'' Roboco giơ ngón tay cái: ``Vậy là thành công
rồi! Lần sau tớ sẽ làm tốt hơn nữa.''

Câu nói khiến mọi người phá lên cười, nhất là khi nhìn thấy vẻ mặt nửa vui nửa ngượng
của Miko. Fubuki, với giọng hài hước, đùa thêm: ``Nếu có lần sau, nhớ chọn ai đó bình
thường hơn làm đồng đội, Miko-chi\~{}''

Nàng Vu nữ bĩu môi, còn nàng Người máy thì đổ mồ hôi hột cười xòa. Không khí tràn
đầy tiếng cười và sự vui vẻ khiến cả nhóm quên đi những khó khăn lúc làm bánh.

\ct

Sau khi nếm thử và thảo luận kỹ càng, Sora, Mel, Mio, Fubuki và Kuon cùng nhau đưa
ra quyết định cuối cùng về giải thưởng.

``Anh có được nhắc rằng là cả tám đứa đã tổ chức một cuộc thi làm bánh mochi, mà đúng
hơn là Ayame-chan bày trò, đúng chứ?'' cậu lên tiếng. Cả tám cô gái đều gật đầu, trong
khi nàng Quỷ sừng thì đang cầm máy quay để ghi lại mọi thứ.

``Được rồi, theo lời của Ayame-chan, thì đội nào thắng sẽ được giải là một phong bao lì
xì, đúng chứ?'' cậu hỏi tiếp. Cả tám cô gái tiếp tục gật đầu.

Lần này, Sora lên tiếng: ``Bọn chị đã thảo luận với nhau khá kỹ càng rồi. Như vậy, cả bốn
đội mọi người đều là những người chiến thắng!''

Cả tám cô gái đều hò reo mừng rỡ, như thể họ là những đứa trẻ sắp nhận được quà. Kaien
mỉm cười, rút ra phong lì xì đỏ tươi, chia đều cho tất cả mọi người: ``Giải thưởng đầu tiên,
lì xì cho tinh thần làm việc chăm chỉ của các con, sẽ được chia đều. Ai cũng xứng đáng
cả!''

Tiếng vỗ tay reo hò vang lên, không khí tràn ngập niềm vui. Tiếp đó, Fubuki lên tiếng:
``Và giờ, giải đặc biệt: một đêm ngủ chung với Kuon-senpai!''

Ngoại trừ cô em gái Ryuuhou, bảy người còn lại tỏ vẻ phấn khích hơn hết. Nhưng khi
mọi người định hỏi về việc đó thì ``nhân vật trung tâm'' đã biến mất từ lúc nào.

Tất cả mọi người, kể cả bốn cô gái ``ban giám khảo'' cũng đều sững sờ trước sự biến mất
khó hiểu của cậu.

\begin{dialogue}
  \speak{Miko} Ể!? Kuon-senpai đâu rồi, ne!?
  \speak{Rushia} *ngại, tức* Đừng hòng chạy trốn, Kuon-senpai! Em đã cố hết sức vì phần thưởng này mà!
  \speak{Noel} Chắc anh ấy sợ chết khiếp rồi. Nhưng mà mình nghĩ chắc anh ấy chưa chạy đi quá xa đâu.
\end{dialogue}

Tất cả mọi người gật đầu lia lịa. Bất chợt, nàng Miêu lên tiếng: ``Mà này\dots\ liệu bảy người chúng ta cùng chia sẻ giải thưởng này thì có được không?''

Bạn của cô vừa giơ tay lên cao, vừa cười lớn: ``Được chứ! Tôi thì chỉ cần một góc là đủ
rồi ha\~{}!''

Nàng Pháp sư chạy tới cô em gái - cũng là bạn đồng hành của cô - tra hỏi: ``Ryuuhou-chan,
anh trai của em chạy ở đâu?'' Ryuuhou nhìn đôi mắt của nàng Pháp sư đang bùng lửa kia,
vội chỉ về cửa ở phía bên trái: ``E-Em nghĩ là Onii-san đi về bên kia ạ\dots''

Thế là, cả bảy người cô gái - Rushia, Noel, Miko, Roboco, Okayu, Korone, và Pekora -
lao vào một cuộc rượt đuổi khắp nhà tìm cho ra cậu con trai. Các cô gái không ngừng gọi
tên Kuon, mỗi người một giọng, tiếng chân dồn dập vang lên khắp hành lang lát gạch, cửa
kéo bằng gỗ liên tục mở ra đóng vào dồn dập.

Kuon, người đang ôm lấy một chiếc gối như vũ khí tự trấn an, len lén trườn qua sau chiếc
tủ chén, rồi lách qua cánh cửa trượt để ra hành lang: ``\textit{Bình tĩnh nào Kuon\dots\ rẽ trái, leo
  thang, trốn sau rèm\dots\ Đừng để họ bắt được!}'' Cậu nín thở khi nghe thấy tiếng dép geta
cộp cộp tiến lại gần, rồi nhanh như chớp bò thấp người chạy ngang qua chiếc bàn kotatsu,
suýt nữa thì vấp phải dây điện. ``\textit{Aish\dots\ Mi-chan mà thấy là xong đời\dots}''

Nàng thỏ Pekora thì tìm ở tầng một ngôi nhà: ``Kuon-senpai\~{}! Đừng có mà trốn, peko!
Bọn em chỉ muốn nói chuyện thôi mà, peko!'' Cô vừa gọi vừa thò đầu xuống gầm ghế sofá,
lôi ra một củ cà rốt làm ``mồi nhử,'' đặt giữa phòng khách: ``Nếu anh ra đây, peko, em sẽ
tặng anh cà rốt đặc sản\dots\ với một vé ngủ chung, peko!'' Tiếng chân cô chạy vòng quanh
chiếu tatami loạt xoạt, kéo cả tấm rèm cửa lên xem phía sau, nhưng chỉ thấy chiếc quạt
đang quay chầm chậm.

Roboco thì tìm trong tầng năm, bên nhà Hijaki: ``Kuon-senpai! Em hứa là sẽ không bật
chế độ siêu tốc nữa đâu! Mau ra đây đi nào!'' Cô chạm ngón tay lên bên tai, bật chế độ
``tìm kiếm thủ công'' để giữ lời hứa, mắt phát ánh xanh nhạt, quét từng phòng bằng nhịp
điệu đều đặn: \textbf{*pi-pi-pi*}. ``Không turbo, không tên lửa, chỉ là robot dịu dàng\dots'' Cô khẽ
nhấc một hộp đồ, nhìn xuống sau kệ sách: ``Hmm\dots\ ở đây chỉ có băng video thôi.''

Miko thì tìm trong căn bếp: ``Kuon-senpai, anh đang trốn ở đâu vậy ne\~{}? Ra đây đi nào!''
Cô lục tủ bếp, vô tình làm rơi một nắp nồi kêu \textbf{*keng!*}. ``Ara ara\dots\ Xin lỗi nha ne\~{}!'' Miko
núp sau cánh cửa tủ lạnh, ghé mắt nhìn dưới gầm bàn, rồi mở tủ gia vị: ``Đây chỉ có mắm
muối thôi ne\~{}!'' Cô cầm chiếc môi gỗ như cây đũa thần, đi khắp bếp, gọi tên cậu bằng giọng
réo rắt.

Noel tìm ở các kho mà họ đã lau dọn từ ngày hôm kia: ``Kuon-senpai! Em sẽ không làm
gì với anh đâu! Em chỉ muốn nói chuyện thôi mà!'' Nàng Kỵ sĩ vén tay áo, mở từng cánh
cửa kho bên cầu thang, nhấc những thùng các-tông to đùng, rồi đặt xuống hết sức cẩn trọng để không làm đổ đồ đã sắp xếp. ``Anh có trốn trong thùng này không đó?''

Rushia thì hớt hải hơn: cô liên tục chạy lên tầng bốn, năm và sáu, lục soát khắp mọi nơi.
Cô không nói một câu gì, nhưng nhìn cách cô tìm cả ngôi nhà như vậy cũng điểu vì sao
cô lại nhanh như vậy.

Cặp đôi Okayu-Korone thì tìm ở bên phòng của Kuon, dẫu vậy họ đang đùa nghịch với
nhau hơn là đang đi tìm cậu con trai. Tréo ngoe thay, đây lại là nơi mà Kuon đang trốn
trong đó, nhưng chưa phát hiện ra\dots

\begin{dialogue}
  \speak{Okayu} Hmm\~{} Không biết Kuon-senpai đi đâu rồi ha\~{}\dots
  \speak{Korone} Ogayu! Tôi đánh được mùi anh ấy kìa!
\end{dialogue}

\dots hoặc là không. Ngay tức khắc, nàng Khuyển sủa lên mấy tiếng inh ỏi, gây chú ý tới con chó Mi nhà cậu.

Thấy cô chó đang lon ton chạy tới, Okayu nói, miệng nở một nụ cười ``nồng thắm'': ``Mi-chan, em hãy gọi những người khác lên đây nha\~{}''

Ngay sau đó, cô chó Mi sủa inh ỏi, chạy khắp nhà, hòa cùng với tiếng sủa của Korone, gây chú ý tới năm cô gái còn lại.

``\textbf{Đâu, đâu!? Kuon-senpai ở đâu!?}'' nàng Pháp sư là người chạy vào nhanh nhất. Theo
sau đó là những thành viên còn lại. Ryuuhou và Ayame theo sau, với cô em gái thì tặc
lưỡi cười mếu cho ông anh trai, còn nàng Quỷ sừng thì cười lên khoái chí.

\begin{dialogue}
  \speak{Ayame} Có vẻ như Onii-san vẫn không thể thoát khỏi họ nhỉ, yo\~{}
  \speak{Ryuuhou} Làm anh phải vất vả rồi, Onii-san\dots\ Tội nghiệp anh thật đó.
\end{dialogue}

Cậu con trai cuối cùng bị tìm thấy đang trốn sau một tấm rèm trong phòng khách. Nhìn
thấy các cô gái tiến đến gần, cậu thở dài, đưa tay lên đầu hàng: ``Được rồi, được rồi! Anh
chịu thua!''

``Bắt được anh rồi nha, senpai!'' cả bảy người chạm nhẹ vào người cậu, như thể họ đang
chơi trò mèo vờn chuột.

Kuon đứng dậy, phủi bụi trên người, rồi chống hông nhìn cả bảy cô nàng, ôm trán:
``Nhưng\dots\ làm thế nào mà tất cả lại đều đòi ngủ chung thế này?''

Cô đồng bước tới, nở một nụ cười ngọt ngào, nhưng lại đầy áp lực: ``Đây là công lao
của chúng em đó, Kuon-senpai. Đừng lo, bọn em sẽ không làm phiền gì đến anh đâu\~{}''

Nhìn nét mặt quyết tâm của các cô gái, Kuon chỉ biết cười trừ. Cuối cùng, để chiều lòng
tất cả, cậu quyết định: ``Thôi được, anh đồng ý. Nhưng chỉ lần này thôi, rõ chưa?''

Tiếng reo hò vang lên như chiến thắng, trong khi A-chan chỉ đứng từ xa, tay ôm trán:
``Đúng là nhà này càng lúc càng náo loạn. Mình nên chuẩn bị tinh thần cho những chuyện
tương tự vào năm sau thôi\dots''

Cô em gái thì lắc đầu ngán ngẩm: ``Đúng là anh ấy không bao giờ yên bình với bọn họ
nhỉ.'' Ayame thì cười một cách ranh mãnh: ``Tết này sẽ có nhiều thứ hay ho lắm đây,
yo\~{}''

Cả bảy cô gái vui vẻ cùng Kuon chuẩn bị một buổi tối đáng nhớ, còn cả nhà thì không
ngừng cười trước tình huống vừa bất ngờ vừa hài hước này.

\ct

Sau bữa trưa với bánh mochi đầy náo loạn và hài hước, cả nhóm chuẩn bị lên đường về
quê ở ngoại thành Kawachi để thăm nhà ông bà nội của Kuon. Vì ông bà nội đã lâu chưa
về thăm quê, Kuon quyết định chia nhóm để tiện việc di chuyển.

Các thành viên còn lại của \hll\ cũng đã tươm tất quần áo truyền thống, do chính họ
mang tới. Ai nấy đều toát ra một vẻ đĩnh đạc và cổ điển, nhưng vẫn không giấu nổi sự
tinh quái, thậm chí có người còn mặc bộ đồ trông\dots\ khá lộ liễu.

\begin{figure}[H]
  \centering
  \begin{subfigure}[t]{0.2\textwidth}
    \centering
    \includegraphics[height=8cm]{../../../../Image/Book?/newyear/haato_model.png}
    \caption{Haato diện một bộ
      furisode đỏ với bộ
      kimono in hoa bên ngoài,
      một chiếc khăn lông, dây
      buộc obi vàng, đi tất tabi
      trắng và geta đen. Cô để
      tóc rất dài với hai bên tóc
      tết và có sợi ahoge, cài
      bởi một chiếc nơ nhỏ và
      một bông hoa trắng.}
  \end{subfigure}
  \quad
  \begin{subfigure}[t]{0.2\textwidth}
    \centering
    \includegraphics[height=8cm]{../../../../Image/Book?/newyear/subaru_model.png}
    \caption{Subaru diện bộ
      kimono in họa tiết uroko,
      áo choàng ren xanh, váy
      hakama xanh cạp cao và
      đôi bốt buộc dây nâu. Cô
      để tóc như thường ngày,
      nhưng được trang trí
      bằng nơ tóc hai màu
      trắng-đỏ cùng 2 bông
      hoa cài tóc.}
  \end{subfigure}
  \quad
  \begin{subfigure}[t]{0.2\textwidth}
    \centering
    \includegraphics[height=8cm]{../../../../Image/Book?/newyear/matsuri_model.png}
    \caption{Matsru mặc bộ vu nữ
      phi truyền thống:
      kimono hở vai màu trắng
      có tay áo rộng tách rời,
      buộc bằng obi cam ở trên
      và obi sọc trắng-vàng ở
      dưới giữ váy ngắn
      hakama màu đỏ. Cô đi
      tất tabi trắng và dép
      zouri. Tóc cô được cắt
      ngắn, với một vệt màu
      xanh duy nhất.}
  \end{subfigure}
  \quad
  \begin{subfigure}[t]{0.2\textwidth}
    \centering
    \includegraphics[height=8cm]{../../../../Image/Book?/newyear/aqua_model.png}
    \caption{Aqua diện bộ hầu gái
      wa: bộ kimono màu
      xanh có xếp nếp mặc bên
      ngoài trang phục thường
      ngày, tạp dề trắng có xếp
      nếp ở eo, tất bobby trắng
      có họa tiết mỏ neo và
      khăn okobo xanh. Tóc ô
      buộc một bên bằng dây
      thừng màu xanh, búi tóc
      lệch sang trái.}
  \end{subfigure}
  \quad
  \centering
  \begin{subfigure}[t]{0.3\textwidth}
    \centering
    \includegraphics[height=8cm]{../../../../Image/Book?/newyear/aki_model.png}
    \caption{Aki đội chiếc mũ nồi xanh, bộ
      kimono không tay sọc dọc xanh, tay
      áo rộng và có cùng họa tiết với
      kimono, ống tay áo bên trong, áo
      choàng đỏ, váy hakama dài màu đỏ,
      tạp dề eo màu trắng có xếp ly và đi bốt
      buộc dây màu nâu. Cô để kiểu tóc bob
      bất đối xứng: phần mái hất lệch, phần
      còn lại được buộc búi cao một bên lai
      kiểu tóc đuôi ngựa lệch, được cài
      bằng hoa kanzashi và nhiều kẹp tóc.}
  \end{subfigure}
  \quad
  \begin{subfigure}[t]{0.3\textwidth}
    \centering
    \includegraphics[height=8cm]{../../../../Image/Book?/newyear/shion_model.png}
    \caption{Shion diện bộ kimono ngắn in họa
      tiết tím với tay áo xếp nếp mặc bên
      ngoài bộ váy mặc định, hai dây obi
      hồng-đen, đi quần tất đen, đai chân
      trái và bốt viền lông đen. Cô có quàng
      khăn lông, giữ bằng trâm cài ngôi sao
      năm cánh. Cô để búi tóc lệch sang
      phải, trang trí bằng hoa kanzashi tím
      tua rua và một dây buộc tóc xanh lá
      hình hai tam giác lồng nhau.}
  \end{subfigure}
  \quad
  \begin{subfigure}[t]{0.3\textwidth}
    \centering
    \includegraphics[height=8cm]{../../../../Image/Book?/newyear/choco_model.png}
    \caption{Choco diện bộ kimono đỏ trễ vai
    mặc bên ngoài một chiếc váy bút chì
    ngắn màu đen có khe ngực lưới xuyên
    thấu và dép okobo đen. Tóc cô tết
    quanh đầu với búi tóc một bên được
    cố định bằng kẹp tóc và hoa kanzashi.
    [DIGITAL}
  \end{subfigure}
  \label{setthree}
\end{figure}

Shion, người đã thành thạo phép chuyển dịch, được giao nhiệm vụ đưa ông bà nội, Mel,
Sora, Ayame, Pekora và Aqua đi trước. Đứng trước khoảng sân, nàng Phù thủy trẻ tuổi
nghiêm túc nói với hai người lớn tuổi nhất trong gia đình: ``Chuẩn bị sẵn sàng đi nhé, các
cụ! Đừng lo, cháu làm việc này cả trăm lần rồi!''

Mel liếc nhìn nàng Phù thủy, giọng đầy nghi ngại: ``Em có chắc là chúng ta sẽ ổn không
vậy? Chị nghe nói em đã từng tự dịch chuyển mình bên trong tường đấy\dots\footnote{Xem {\flaregothic ÉPISODE 34}, quyển số Ba.}''

``Cũng đừng quên là do ma thuật của chị mà em bị kẹt lại ở trong bức tường đó, peko,''
nàng Thỏ lườm Shion, tỏ vẻ không vui khi cô nhớ lại khoảnh khắc đáng xấu hổ đó.

Aqua vẫn đang loay hoay buộc đồ đạc, hoảng hốt quay lại: ``Ể!? Lần trước là thế nào
vậy? Shion-chan, bà chắc chứ?''

Nghe cô bạn đồng thế hệ nói vậy, cô phồng má, giậm chân: ``Mọi người vẫn không tin tôi
à!? Lần này tôi đã cải tiến thuật pháp rồi!''

``\textit{Liệu chị ấy có đáng tin không đây\dots}'' Ryuuhou thở dài, dù không hiểu chuyện gì đang
xảy ra hay đang nói gì. Kuon bật cười, vỗ vai nàng Phù thủy đang nóng máu kia: ``Anh
tin em, Shion-chan. Nhưng mà có sai sót thì đưa họ về ngay lập tức, hiểu chứ?''

Nghe cậu nói vậy, nàng ``Tỏi'' càng tức điên hơn: ``Rồi, rồi, em biết rồi! Đừng có mà chỉ
đạo em nữa!''

``Anh yên tâm, bọn em lo liệu được!'' Sora nói, nở một nụ cười thắm. Ayame cũng không
kém cạnh: ``Bọn em luôn có chiêu hết rồi, yo!''

Shion hắng giọng một tiếng, rồi đọc một câu thần chú. Ngay sau đó, ánh sáng tím lóe lên,
rồi lóa sang một thứ ánh sáng trắng bao trùm cả không gian rộng lớn, và rồi\dots\ nhóm đầu
tiên biến mất vào không khí.

Kaien đứng bên cạnh, lắc đầu: ``Cứ như thế này thì mẹ không ngạc nhiên nếu có ai đó lại
lạc đường\dots'' Cậu con trai chỉ cười xòa: ``Con cũng không tin em ấy đâu. Nhưng mà con
tin là nếu có vấn đề gì thì Mel-chan và Sora-chan sẽ lo được mà. Con cũng đưa cho họ
địa chỉ rồi mà.''

Các thành viên còn lại chia ra thành ba nhóm. Gia đình Kuon di chuyển bằng hai xe máy,
trong khi nhóm Fubuki, Mio và Subaru đi trên chiếc xe cảnh sát siêu tốc của Cáo Trắng.
Những người còn lại chen chúc trong chiếc xe hơi ``Lamb*ghini\footnote{Sẽ được giới thiệu ở quyển số Bốn.}'' của Okayu.

Ngay khi vừa lên xe, Fubuki liền đạp ga mạnh mẽ, khiến Subaru hét lên: ``Oi, oi! Chúng
ta đi thăm quê, không phải đua xe đâu!''

Trong khi đó, tài xế của chiếc xe đeo kính mát, cười lớn: ``Đi nhanh chút cho có cảm giác
mạnh mà! Em nên cảm ơn chị vì tiết kiệm thời gian đấy, Subaru-chan!''

Mio ngồi ghế sau, chỉ biết thở dài mà giữ chặt tay vịn: ``Họ sẽ phạt chúng ta mất nếu vượt
tốc độ á. Lái xe như bình thường đi, Fubuki ạ\dots''

Rushia, Noel, Matsuri, và Miko, những người chỉ im bặt vì tốc độ điên rồ mà chiếc xe của
cảnh sát đang phóng đi.

Cùng lúc đó, xe của Okayu cũng chẳng khá hơn. Korone, ngồi ở ghế lái phụ, hò hét:
``Nhanh lên Ogayu, họ sắp vượt mặt rồi!''

Nàng Miêu bình thản nhấn ga thêm: ``Bình tĩnh, bình tĩnh. Tôi cũng không muốn bị bắt
đâu ha\~{}''

Các thành viên còn lại như Aki, Choco, Roboco, Haato, A-chan thì không dám mở miệng
vì họ đi quá nhanh.

Xe của Kuon-Ryujin và Kaien-Ryuuhou, lái xe máy, theo sau cả hai xe ô tô. Nhìn cảnh
tượng trước mắt - hai chiếc xe đua nhau trên đường - cậu Tổng chỉ biết thở dài.

``{\color{red} Các bạn của anh đúng là điên rồ thật đấy,}'' cô em gái của cậu đổ mồ hôi hột khi trông
thấy hai chiếc xe đang vừa phóng như bay, vừa kèn cựa nhau, tới mức một chiếc xe máy
của người tham gia giao thông còn suýt ngã ra đường.

``{\color{red} Đồng nghiệp. Nhưng đúng, đúng là điên thật. Họ tưởng rằng đây là một cuộc đua hay sao\dots}'' cậu sửa lại, ngán ngẩm dõi theo ánh mắt của cô em gái.

Ryujin, từ chiếc xe máy bên cạnh, bật cười: ``{\color{red} Cũng may là không có công an}.'' Kaien,
người đang cầm lái chiếc xe kia, khẽ nhíu mày: ``{\color{red} Anh đừng có vớ vẩn. Hy vọng là họ sẽ
  đi bình an\dots}''

``{\color{red} Con nghĩ mẹ đang lo bò trắng răng đấy. Họ chắc chắn sẽ ổn thôi, con biết mà,}'' Kuon
đáp, dẫu cho lời nói của cậu có phần ngập ngừng. Cô em gái đi xe bên cạnh, cũng chỉ tặc
lưỡi: ``{\color{red} \dots Anh đúng là lạc quan tếu thật.}''

Cả bốn người trong gia đình chỉ biết lắc đầu, tiếp tục lái xe, để mặc hai chiếc ô tô phóng
đi phía trước như một cơn lốc.

\ct

Hai chiếc xe do Fubuki và Okayu lái phóng như tên lửa trên con đường quốc lộ. Cảnh
tượng không khác gì một màn chơi A(s)phalt sống động: tiếng động cơ gầm rú, bánh xe
nghiến trên mặt đường, và cả những tiếng la hét từ hành khách phía sau.

Fubuki, mắt sáng rực như đã vào ``chế độ game thủ'': ``Cứ thế này tôi sẽ về nhất cho xem!''
Đối thủ của cô, Okayu, tay vững vàng trên vô lăng, bật cười: ``Không, không, là tôi chứ\~{}''

Mio cố giữ không ngã khỏi ghế sau, hét lớn: ``\textbf{Đừng có tự biến mình thành hung thần
  đường phố chứ!!}''

Korone phấn khích, đến mức không ngừng cổ vũ: ``Ogayu cố lên, cố lên nào\~{}'' dường
như ngó lơ những hành khách đang tái mét mặt mày ở phía sau.

Noel ngồi ở phía sau, nắm chặt ghế trước, hét lớn: ``\textbf{Cố lên, Fubuki-senpai!!}''

Subaru cố níu chặt dây an toàn, la to: ``\textbf{Đừng có cổ vũ người ta chứ! Đường phố này có
  giống như ở thành phố Điện tử đâu!}''

Rushia cũng không khá hơn, khi cô nàng đang co rúm trong ghế, run rẩy nói: ``Đ-Đáng
sợ thật đó-nanodesu\dots''

Hai chiếc xe của Fubuki và Okayu tiếp tục phóng với tốc độ kinh hoàng ``đấu trường đua
xe.'' Những chiếc xe khác vội tấp sát lề đường để tránh hai chiếc xe hung thần lao qua,
trong khi các hành khách không ngừng hét lên hoảng loạn.

Nàng Sói đen vẫn đang bám chặt vào ghế trước trong xe của nàng Cáo trắng, hét lớn:
``\textbf{Xin bà đó, Fubuki! Ta đang ở ngoài đường đó, không phải đang quay phim hành
  động!}'' Đằng sau cô, Subaru ôm đầu: ``Em không muốn làm diễn viên đóng thế đâu!''

Nhưng thay vì nghe lời bạn của mình, Fubuki lại nhấn ga mạnh hơn, tỏ vẻ rất phấn khích:
``Mio\~{} tin tôi đi, đây sẽ là một chuyến đi điên rồ đấy!''

Đúng lúc đó, Fubuki trông thấy một chiếc xe tải lớn đang chạy phía trước, trên đó có một
ván dốc chở hàng. Mắt cô sáng lên: ``Nhìn nè, Mio! Một cú xoay vòng tuyệt hảo đang
chờ chúng ta!''

Noel thấy vậy, nổi hứng hơn hẳn: ``Tới luôn đi, Fubuki-senpai!!'' mặc cho những lời la ó
can ngăn của những người còn lại trên xe.

Nàng Cáo trắng đánh lái và tăng tốc, lao thăng lên ván dốc xe tải. hiếc xe lập tức nhấc
bổng lên không trung, xoay một vòng tròn hoàn hảo trên không\dots

\begin{dialogue}
  \speak{Miko} Cái gì vậy ne\dots?
  \speak{Mio} \textbf{Fubuki!!}
  \speak{Fubuki} Woo-hoo!
  \speak{Noel} Ta đang xoay vòng trên cao!
\end{dialogue}

\dots trước khi đáp xuống mặt đường với một cú nảy mạnh, khiến những người đi đường
thêm một phen khiếp đảm.

Mio lúc này tái mét, cả người cô mềm nhũn ra: ``Tôi sẽ không đi chung xe với bà một lần
nào nữa đâu!'' Rushia và Miko ôm chặt lấy nhau, gần như khóc:

\begin{dialogue}
  \speak{Miko} T-Thật là đáng sợ quá, ne\dots\ Chị tè ra quần mất thôi!
  \speak{Rushia} B-Bình tĩnh lại đi nào, Miko-chan! Fubuki-senpai, dừng lại đi!
\end{dialogue}

Ở phía sau, xe của Okayu cũng không chịu kém cạnh. Korone, mắt sáng rực như đèn pha,
chỉ tay về phía một chiếc ô tô đỗ bên lề đường: ``Ogayu! Chúng ta phải làm gì đó thật ấn
tượng hơn!''

Nàng Miêu liếc nhanh, nhoẻn miệng: ``Bà đang có bài gì sao, Koro-san?'' Nàng Khuyển
hét lớn, giơ nắm đấm lên cao không tỏ vẻ do dự gì: ``\textbf{Lộn vòng đi nào!}''

A-chan vội lên tiếng: ``N-Này! Các cô tính làm cái trò gì vậy!?'' nhưng dường như câu
nói của nàng Đeo kính không thể đến tai của cặp đôi ``quỷ sứ'' này.

``Vậy thì lên nào!'' Okayu cười, đánh mạnh vô lăng. Xe lập tức xoay một vòng tròn trên
đường nhựa, trước khi bất ngờ đâm nhẹ vào mép vỉa hè, nhấc bánh sau lên, và\dots\ chiếc xe
lao vào không trung như một pha nhào lộn trong phim hành động của xứ cờ hoa.

Những hành khách trên chiếc xe này cũng la hét lớn không kém gì xe đi đằng trước họ:

\begin{dialogue}
  \speak{A-chan} Inugami-san! Nekomata-san! Các cô sẽ bị trừ lương vì chuyện này\dots!
  \speak{Aki} Eeek!
  \speak{Choco} Hai em đang làm cái trò gì thế hả!?
  \speak{Haato} \textbf{Xoay vòng nào, cặp đôi \textit{GAMERS} ơi!}
\end{dialogue}

Chiếc xe đáp xuống với một cú nảy mạnh, tông hạ luôn chiếc ô tô đỗ bên lề, khiến nó bật
tung ra một góc.

A-chan ngỡ ngàng, tái mét sau cú lộn vòng đó: ``T-Tôi không thể tin là hai cô vừa làm thế
thật\dots''

Choco cũng kinh hoàng không kém: ``Mấy đứa làm rơi hết đồ phấn trên mặt cô rồi đó!!''
Aki đáp lại đầy ngỡ ngàng: ``Đó là điều mà cô quan tâm sao!?''

Haato, ngồi cạnh A-chan, nhìn cảnh tượng qua cửa sổ, la lớn: ``\textbf{Âu mài gót!! Chị phải
  ghi lại cảnh này làm nội dung ngay!}'' và hú hét trong lúc cửa sổ vẫn đang mở.

Phía sau, Roboco cười vui vẻ, như thể không có chuyện gì đáng sợ xảy ra: ``Tuyệt quá,
Okayu-chan! Lần sau làm thêm cú nhào lộn nào khác đi!''

Những người qua đường thấy hai chiếc xe đang lao nhanh về phía trước liền bỏ chạy tán
loạn, một người thậm chí còn suýt bị một trong hai chiếc xe chẹt, nếu không vội bỏ chiếc
xe đi.

``{\color{red} C-Có phải họ vừa biến đường phố này thành một trường phim hành động không vậy!?}''
tất cả bọn họ đều kinh hãi trước hai chiếc xe vừa có ``màn trình diễn'' điên rồ đến thế.

Hai chiếc xe lại tiếp tục tăng tốc, chuẩn bị cho thử thách tiếp theo phía trước\dots

\ct

Okayu và Fubuki vẫn đang tiếp tục kèn cựa nhau, vừa đi vừa lạng lách qua các xe cộ khác,
gây náo loạn thêm tại một cung đường nữa.

Họ ngày càng bám sát nhau hơn, thậm chí còn tông đổ cả dải phân cách để chạy sang làn
bên cạnh, càng khiến cho nhiều người đi đường ở hướng ngược chiều thêm khiếp đảm.

Bất ngờ, nàng Dị giới bất ngờ cất giọng lo lắng: ``Ờm, mọi người\dots?'' Tất cả mọi người
ở cả hai xe đồng loạt vang lên: ``Sao!? / Cái gì!?''

``Phía trước có tàu hỏa kìa\dots''

Tất cả đều quay đầu nhìn về phía trước. Một đoàn tàu hỏa khổng lồ đang chậm rãi tiến
lại gần, và hàng rào chắn đang từ từ hạ xuống.

``Mau dừng xe đi chứ, Okayu-sama! Đừng có nghĩ đến chuyện dại dột!'' nàng Y tá hoảng
hốt. Nàng Pháp sư ở xe của Fubuki cũng lo lắng không kém: ``Em là Cô đồng nhưng em
cũng không muốn chết đâu!!''

Trong khoảnh khắc đáng lẽ phải dừng lại, những ``quái xễ'' lại đồng loạt nảy ra một ý
tưởng ``táo bạo''. Dường như họ chẳng đếm xỉa gì tới những tiếng la hét thảm thiết của
các hành khách phía sau.

``\textbf{Ta cùng vượt đường ray thôi!!}''

Noel và Roboco, cặp đôi ngốc nghếch nhưng hăng máu (vì không sợ chết) nhất, reo hò:
``Vượt rào! Vượt rào!'' Haato thậm chí còn hét lớn: ``\textbf{Nhanh nữa lên nào!!}''

Miko lúc này đã sợ đến mức ngất xỉu, khiến cho Rushia ngồi cạnh: ``Miko-chan!? Mikochan!!
Đừng chết mà!'' Nàng Vịt cũng không kém cạnh: ``Đây có phải là GTA đâu!!''

Nàng Song tính tiếp tục hú hét, cổ vũ cho cả hai người: ``\textbf{Cố lên nào mọi người!}'' Lập
tức, Subaru không kiềm được cơn giận: ``Lại còn cả chị nữa!! Ngậm miệng lại cho em
đi!!''

Hai chiếc xe lao tới với tốc độ không tưởng, bất chấp hàng rào chắn tàu hỏa đang hạ xuống
hoàn toàn. Khi đoàn tàu chỉ còn cách vài mét, cả hai xe đồng loạt tông sầm vào hàng rào
chắn, làm nó bật tung lên.

Toàn bộ những người bên trong xe lại tiếp tục la hét. ``\textbf{Chúng ta sẽ chết mất! Chúng ta
  sẽ chết mất thôi!!}'' một vài người trong số họ thất thanh, trong khi một số người thì kích
động vui sướng tiếp tục cổ vũ.

Trong tích tắc, hai xe vượt qua đường ray ngay trước khi đoàn tàu lao qua. Tiếng còi tàu
rú vang, làm chấn động cả khu vực. Những người tham gia giao thông nhìn cảnh tượng
đó trố mắt ra ngỡ ngàng, không thể tin vào những gì họ đang được chứng kiến.

Hai quái xế sau khi đã lập ``kỳ tích'' trong sự khó hiểu, kinh hãi và sửng sốt của mọi người, họ vui vẻ đập tay nhau qua cửa sổ xe:

\begin{dialogue}
  \speak{Fubuki} Whoo! Quá tuyệt vời!
  \speak{Okayu} Lần sau ta lại làm thế nha\~{}
\end{dialogue}

Trong khi đó, các hành khách trên hai xe, ai nấy đều sợ đến tái mặt, tiếng hét vẫn còn
vang vọng. Miko giờ đây đã bất tỉnh, tới mức độ hồn của cô lìa khỏi xác. Rushia thấy
vậy, hoảng hốt: ``\textbf{Em muốn được xuống xe!!}''

Nàng Tomboy run rẩy, tay chỉ về phía trước: ``L-Làm ơn\dots\ đừng đua nữa\dots!'' Nàng Sói
đen thì ôm ngực, thở dốc: ``Tôi thề là sẽ không bao giờ ngồi xe của các cậu nữa!''

Bỏ mặc những lời than phiền, hai chiếc xe lại tiếp tục phóng đi, để lại đằng sau một khung
cảnh hỗn loạn ngổn ngang với những người qua đường đang ngơ ngác nhìn theo, tự vấn
không biết những chuyện quái quỷ gì họ vừa thấy.

\ct

Hai chiếc xe tiếp tục phóng trên đường cao tốc, xung quanh là những cánh đồng xanh
ngát kéo dài đến tận chân trời, thi thoảng họ còn bắt gặp một vài ngôi nhà mọc ra khá thưa
thớt.

``Có vẻ\dots\ mọi thứ trở lại bình thường rồi,'' nàng Đeo kính thở phào nhẹ nhõm, ``ở đây
cũng gần như không có gì để cho mấy đứa này quậy nữa.''

Như những gì cô nói, đoạn đường cao tốc này khá mới, lại gần như không có xe cộ nào,
đường lại khá êm, nên hầu như không có bất cứ thứ gì để cho hai ``quái xế'' kia thỏa sức
tung hoành được nữa.

``Ẻ\~{} Chán òm à\~{}'' Haato chui vào trong xe, tỏ vẻ hơi thất vọng. Roboco cũng chỉ lầm lũi
ngồi đó, mắt giương về cửa sổ. Những người còn lại thì trông như thể lấy lại được sự
sống.

Khi không khí tưởng chừng đang yên bình, một bóng đen bất ngờ xuất hiện từ trên trời,
đâm thẳng xuống mui chiếc xe của Okayu và Korone, tạo nên một tiếng *rầm!* cực lớn.

Choco giật mình thảng thốt: ``Cái quái gì vừa rơi xuống thế!?'' Tiếng động mạnh cũng
khiến cho mọi người trên xe chú ý.

Nàng Miêu vẫn giữ nguyên vẻ bình thản như thường lệ, liếc qua gương chiếu hậu rồi nhún
vai: ``Ồ, có vẻ như thứ gì đó đâm trúng xe rồi.'' Nàng Khuyển cười trừ: ``Tiếc ghê\~{} Chiếc
xe của bà bị móp mất rồi.''

Những người còn lại trên xe thì hoàn toàn không thể bình tĩnh như cặp đôi ``bình chân
như vại'' ấy.

A-chan ôm chặt lấy trên ghế, cố gắng lấy lại bình tĩnh: ``Đây có phải là lúc tiếc của đâu!''
Aki bật dậy: ``Chúng ta bị tấn công rồi! Địch đã phát hiện ra chúng ta!''

Haato nghe vậy, đột nhiên tỏ vẻ phấn khích: ``Địch ư!? Đâu, địch đâu!? Chị đây sẽ cho
chúng nó một trận!''

Nàng Y tá thốt lên: ``Làm gì có chuyện đó! Aki-sama, Haato-sama, hai người bình tĩnh
lại-!'' Bất ngờ, chiếc xe mất kiểm soát, chao đảo mạnh và lao thẳng xuống cánh đồng phía
dưới. Mọi người hét toáng lên khi chiếc xe vừa cày nát một mảng ruộng lớn, vừa lao đi
trong tiếng động chát chúa của kim loại và đất đá.

``Chuyện gì đang xảy ra vậy!?'' nàng Vịt hốt hoảng, nhìn đất đá bắn tùm lên cửa kính xe,
nhưng chiếc xe thậm chí còn không có ý định dừng lại.

``Tại sao chúng ta lại ở dưới ruộng chứ!?'' Choco hoang mang. Thậm chí, cô còn hét lớn
khi trông thấy những con đỉa đang bám sát vào cửa sổ.

Trong khi mọi người vẫn đang hoảng loạn, Aki vươn tay giành lấy vô lăng từ tay Okayu:
``Để đó cho chị lái thử nào!''

Không ai kịp phản ứng trước khi Aki bắt đầu điều khiển chiếc xe. Nhưng thay vì ổn định
lại, chiếc xe đột nhiên rùng lên, nhấc bổng khỏi mặt đất và\dots\ bay lên trời!

Nhìn chiếc xe không còn chạm trên mặt đất nữa, nàng Vịt hét lớn, mặt cắt không còn một
giọt máu: ``Aki-senpai, chị vừa làm cái trò gì vậy!? Tại sao nó lại bay chứ!?''

Nàng Song tính thì lại vui thích hơn: cô nàng trèo hẳn ra ngoài xe, đứng cân bằng trên
nóc của chiếc xe: ``Hay quá! Chúng ta như thể là UFO này!!''

Nàng Y tá thì nhắm chặt mắt lại, dường như vẫn đang sợ những con đỉa bám trên cửa kính
hơn là sợ độ cao: ``Mau cho xuống đi! Cô không muốn tham gia trò này nữa!!''

Nàng Người máy vươn người ra ngoài, trái ngược hoàn toàn với vẻ u sầu ban nãy: ``Chị
không ngờ chiếc xe có thể bay được đấy!''

Okayu dường như khá bàng quan với phản ứng với mọi người,thản nhiên giải thích: ``Xe
của chị là Lambo*ghini. Đôi lúc nó thích làm điều kỳ quặc á\~{}''

Nàng Chó vỗ tay động viên: ``Aki-senpai giỏi thật nha\~{}'' Về phía nàng Dị giới, dù đang
bối rối, cố gắng kiểm soát tình hình: ``Ừm, làm sao để nó hạ cánh đây nhỉ?''

Chiếc xe cứ thế bay qua những ngôi làng nhỏ, khiến người dân bên dưới kinh hãi ngước
nhìn. Sau một hồi loạng choạng, nó rơi thẳng xuống\dots\ cổng làng!

Chiếc xe đáp xuống đất với một cú nảy mạnh, tạo nên một cơn bụi mù bao phủ khắp khu
vực. Người dân trong làng hét toáng lên, chạy khắp bốn phương tám phía vì tưởng có
động đất. Trong khi đó, chiếc xe tiếp tục lao về phía trước mà Aki không thể dừng lại.

Nàng Vịt hét lớn: ``Chị không biết phanh à!?'' Nàng Dị giới nghe vậy, luống cuống tìm
đạp phanh: ``Phanh á!? Phanh, phanh, phanh\dots\ đây rồi!''

Ngay khi Aki đạp mạnh bàn phanh, chiếc xe từ từ chạy chậm lại, mọi người trên xe bị hất
về phía trước. Chiếc xe cào thêm một đoạn đường làng khá dài nữa, để lại một vệt đường
bê tông bị cày nát đến mức chẻ đôi.

Chiếc xe của nhóm Okayu-Korone cuối cùng cũng dừng lại, nhưng không phải ở bất kỳ
nơi nào bình thường. Nó đỗ thẳng trước sân nhà ông bà nội của Kuon, nơi Aqua đang
cặm cụi quét lá và Shion vừa lau dọn xong. Tiếng phanh kít đến chói tai làm nàng Hầu
gái giật mình, thét lên một tiếng trước khi bỏ cả chổi chạy thẳng vào trong nhà.

Roboco xuống xe, vẫn cười như thường lệ, vui vẻ thốt lên: ``Đây đúng là trải nghiệm tuyệt
vời!'' A-chan thì loạng choạng bước xuống xe, tìm đến một chỗ nào đó để nôn ra, mệt
mỏi thốt: ``Tôi thề là tôi không còn sức để cho các cô cầm lái nữa đâu\dots''

Nàng Phù thủy vỗ tay chầm chậm, nhìn chiếc xe của cặp đôi Chó-Mèo đang chậm dần,
không giấu nổi nụ cười nhếch mép: ``Ồ, về đích nhanh thật đấy. Nhưng tiếc là vẫn chậm
hơn bố con anh quản lý của chúng ta rồi.''

Nàng Vịt nghe vậy, ngỡ ngàng kêu lên: ``Ể!? Cái gì chứ!? Chúng tôi đã vượt tàu hỏa và
cả bay trên trời nhờ ơn hai con quái xế này cơ mà!'' tay chỉ về cặp đôi Chó-Mèo vẫn đang
giả vờ ngây thơ.

Đúng lúc đó, Kuon từ trong nhà bước ra, khoanh tay nhìn mọi người với vẻ mặt đắc ý:
``Tất cả là nhờ vào power-up thôi.'' Nàng Dị giới bất ngờ khi nhìn thấy cậu Tổng đang đứng sừng sững trước mặt, tò mò hỏi: ``Ý-Ý anh là sao vậy? Anh làm cách nào mà vượt
qua chúng em được chứ?''

Cậu con trai mỉm cười bí hiểm: ``À thì\dots\ nói thế nào nhỉ?''

\begin{center}
  \uline{5 phút trước.}
\end{center}

Gia đình Kuon đang lái xe một cách thong thả trên con đường làng. Tiếng gió vi vu hòa
cùng cảnh đồng quê yên bình khiến không khí vô cùng dễ chịu. Cả bốn người bọn họ vừa
đi, vừa cảm nhận được khí trời đang thổi vào người họ, vừa trò chuyện với nhau vui vẻ
giống như trong gia đình bình thường.

Bỗng nhiên, họ nhìn thấy năm khối lập phương có dấu hỏi đang xoay vòng, lơ lửng giữa
đường. Ryujin nhíu mày nhìn những khối lập phương kỳ lạ: ``{\color{red} Cái gì thế này? Đồ để trang trí xe máy à?}''

Ryuuhou ở xe bên kia, hớn hở nói: ``{\color{red}Cái này con không chắc\dots\ Nhưng chúng trông giống
  power-up trong Ma**o Kart đấy bố!}''

Kuon nghiêng đầu ngạc nhiên: ``{\color{red}Em nói như thể đó là trò chơi vậy\dots\ Mà tại sao nó lại ở
  đây được chứ? Đây là Kawachi, chứ có phải thành phố Điện tử-}''

Tuy nhiên, vì mải quan sát, Ryujin đã vô tình đâm thẳng vào một trong những khối lập
phương làm nó vỡ tung ra. Ngay lập tức, một vòng quay xuất hiện bên cạnh họ, xoay tít
như đang chọn giải thưởng.

``{\color{red}Cái quái gì vừa diễn ra vậy?}'' người bố ngạc nhiên khi vòng quay đang xoay một cách
chóng mặt. Chưa để ai nói gì thêm, vòng quay dừng lại, hiện lên biểu tượng hai mũi tên
màu cam chỉ hướng lên trên.

Ông chớp mắt khó hiểu: ``{\color{red}Đây là biểu tượng gì đây?}'' Cậu Tổng ngồi đằng sau, cầm lấy
hai mũi tên trên và gắn thử vào chiếc yên xe máy ở thùng sau xe: ``{\color{red}Để con thử xem nào\dots}''

Ngay khi mũi tên vừa chạm vào yên xe, chiếc xe máy rít lên một tiếng và\dots\ phóng vọt
lên trời như tên lửa! Người mẹ cầm lái trên xe ở bên cạnh, hoảng hốt hét lên: ``{\color{red}\textbf{Chuyện
    gì thế này!? Xe biết bay ư!?}}''

Cô em gái ngửa mặt lên trời, chỉ vào chiếc xe của hai bố con đang bay cao: ``{\color{red}Onii-san
  đang bay lên trời luôn rồi! Trông như siêu nhân vậy!}''

Kuon dù đang ôm chặt bố, vẫn có thể vọng lại: ``{\color{red}Đừng có cười anh chứ!}'' Đến Ryujin
cũng hoang mang không kém: ``{\color{red}Chúng ta tại sao lại ở trên trời vậy!?}''

Không còn cách nào khác, xe máy của Kaien đành phải đuổi theo, hy vọng tìm cách đưa
Ryujin và Kuon xuống. Nhưng chẳng may, họ lại đâm phải một khối lập phương khác,
và vòng quay lần này dừng lại ở biểu tượng\dots\ lỗ xoáy.
Người mẹ vừa nhìn thấy thứ hiện ra kết quả trước mắt, bàng hoàng hỏi: ``{\color{red}Cái gì đây? Lỗ
  xoáy à? Nó có nghĩa là gì?}''

Ryuuhou tò mò đáp: ``{\color{red} Con cũng không rõ nữa\dots\ Nhưng con thử xem sao!}'' Cô em gái vừa tra biểu tượng lỗ xoáy vào một phần xe, cả hai mẹ con bất ngờ cảm thấy như bị cuốn vào một lực hút mạnh. Trong chớp mắt, họ biến mất khỏi con đường làng và\dots\ xuất hiện
ngay trên nóc xe cảnh sát của Fubuki, vốn đang chạy băng băng trên đường cao tốc.

Trong khi đó, nàng Cáo trắng giật mình khi thấy có một vật thể đột ngột ``rơi'' xuống nóc
xe mình, hét qua bộ đàm: ``Chuyện gì vậy!? Xe của tôi bị đè rồi! Có ai đó liên lạc với
trung tâm không!?''

Mio ở bên cạnh thét lớn: ``Chuyện hú gì đang xảy ra trên kia vậy!?''
Miko thì vẫn chưa tỉnh lại, Rushia thì hốt hoảng, cố lay tiền bối của mình dậy: ``Miko-chan! Miko-chan! Dậy đi, chị ơi! Chúng ta bị tấn công rồi!''

Matsuri, người vẫn đang im hơi lặng tiếng tới giờ, không biết cô nên im lặng hay nên
phấn khích với tình cảnh hiện tại.

Trong khi đó, hai mẹ con nhà Hikawa vẫn đang định hình lại mọi thứ. Kaien ngỡ ngàng:
``{\color{red}Chúng ta đang ở đâu vậy!?}''

Ryuuhou nhún vai: ``{\color{red}Con đâu có biết đâu mẹ.}'' Hai mẹ con bỗng nghe thấy tiếng động từ
phía trên. Họ phát hiện ra vật thể đã đâm trúng nhóm Okayu-Korone trước đó\dots\ chính là
xe máy mà Ryujin và Kuon đang cầm lái trên đó!

Cô em gái tỏ vẻ ngạc nhiên: ``{\color{red}Hai bố con anh ấy đáp xuông như siêu anh hùng đâu!}''
Trong khi đó, cậu Tổng vẫn đang bám chặt vào bố mình, la lớn: ``{\color{red} \textbf{Có cách nào để dừng
    chiếc xe này vậy!?}}''

Ngay lúc đó, chiếc xe máy lại phát động chế độ tăng tốc, lao thẳng về phía trước, bỏ lại
gia đình và cả xe của Fubuki trong mớ hỗn loạn.

Chiếc xe máy của bố con cuối cùng cũng hạ cánh một cácách ngoạn mục ngay trước cửa
nhà ông bà nội của Kuon.

Ryujin, vẫn thở hồng hộc như vừa đi chuyến tàu siêu tốc, nhìn con trai và chiếc xe với vẻ
mặt không thể tin nổi: ``{\color{red}C-Cứ như thể ta vừa đi máy bay về vậy\dots\ Lần sau đừng có thử
  mấy thứ linh tinh nữa, nghe chưa!?}''

Cậu Tổng lắc người, than thở: ``{\color{red}Con đâu có biết nó sẽ thực sự hoạt động thật đâu chứ\dots}''

``{\color{red}Thôi, bố con mình dắt xe vào đi, để xem ông bà nội thế nào rồi,}'' Ryujin nói, vân nhìn cậu
con trai với vẻ trách móc, nhưng sâu trong thâm tâm, đây là lần đầu tiên ông có chuyến
đi kỳ quặc và nhiều cảm xúc đến vậy.

\begin{center}
  \uline{Hiện tại.}
\end{center}

``Và đó là lý do đấy,'' cậu con trai kết thúc câu chuyện với giọng đều đều nhưng không
giấu được vẻ châm biếm. ``Anh có cảm giác mấy đứa đang biến cả thành phố thành một
\textit{hologra} trường kỳ ấy\dots ''

Câu nói vừa dứt, cả nhóm phá lên cười, dù rõ ràng ai nấy đều đang thấm mệt sau chuyến
hành trình ``ly kỳ'' đến khó tin.

Nàng Hầu gái lấp ló trong nhà, rên rỉ: ``Tôi không biết chuyện gì đang xảy ra nữa\dots''

Nhưng chưa kịp định thần, tiếng động cơ rít lên ầm ĩ từ xa báo hiệu một sự cố khác.
Chiếc xe ô tô của nhóm Fubuki-Mio lao tới với tốc độ không kiểm soát, đâm thẳng vào xe
của Okayu, kéo cả hai xuống ruộng bên cạnh. Cảnh tượng diễn ra quá nhanh khiến mọi
người đứng ngoài không kịp trở tay.

Chiếc xe máy của Kaien và Ryuuhou, may mắn hơn, nhẹ nhàng chạm đất mà không gặp
bất kỳ vấn đề nào. Cô em gái xuống xe, vừa hoảng hồn vừa vui sướng: ``Xe của hai người
họ trông vui thật đấy!''

Cậu Tổng khoanh tay đứng nhìn hai chiếc xe giờ đây nhem nhuốc bùn đất, chép miệng:
``\dots Anh nghĩ là mấy đứa chạy quá ẩu đấy.''

Mọi ánh mắt đều đổ dồn vào chiếc xe nằm im dưới ruộng. Haato, lúc này vừa ngồi dậy
vừa la hét đầy uất ức: ``Này, tại sao cứ phải là tôi chứ!?''

Rushia, lập lòe trèo ra khỏi xe, đáp tỉnh bơ đây bất ngờ: ``Do chị xui thôi, Haato-senpai
ạ.'' Cô nhìn xung quanh, ú ớ: ``À-rế, Miko-chan đâu rồi!?''

Họ nhìn thấy nàng Vu nữ đang cắm đầu xuống bờ ruộng, bùn nước thi thoảng còn nổi lên
những bong bóng, cùng với bộ kimono bị lấm lem màu nâu của bùn.

Xung quanh cô là những con cá rô đang vùng vẫy, đòi được trèo lên hẳn\dots\ quần lót trắng.
Cứ một con cá nhảy trúng, cả thân người của cô run lên cầm cập, bọt khí bùn lại càng nổi
lên rõ hơn.

``Không thể ngờ tới luôn á, Miko-senpai thích ăn cá đến mức đó luôn sao?'' nàng Pháp sư
dường như vẫn còn quá sợ, tới mức thốt ra một câu mà chính cô cũng không hiểu tại sao
lại nói vậy. Câu nói của nàng pháp sư nhỏ bé khiến mọi người phá lên cười lần nữa, còn
Miko thì bây giờ bắt đầu vùng vẫy thoát khỏi bờ ruộng.

Trong khi đó, Mio, vẫn đang ngồi trong xe, ôm đầu rên rỉ: ``Chóng mặt quá\dots'' rồi thậm
chí ``gọi tên em Huệ'' (nói trắng ra là nôn) với khuôn mặt xanh xao. Tiếng gọi đột ngột và
biểu cảm mệt mỏi của nàng sói đen khiến cả nhóm không nhịn được cười. Một lần nữa,
không khí trở lại nhẹ nhàng hơn với những tràng cười vang.

\ct

Giữa lúc cả nhóm đang tìm cách giải cứu hai chiếc xe tội nghiệp, một giọng nói trong trẻo
vang lên từ cổng: ``Yoo-hoo, mời mọi người vào nhà nào!'' Họ nhìn thấy nàng Quỷ sừng
đang đứng cười tươi, tay cầm chiếc máy quay mà cô mượn được từ A-chan.

``Ồ? Có vụ gì hay ho kia?'' cô hỏi, tự mày mò tiến gần đến ruộng, nơi mà mọi người đang
cố lôi hai chiếc xe lên mặt đất. Nàng còn nhìn thấy một vài người đang nằm trên vũng
bùn, với nàng Hiệp sĩ và nàng Người máy đang bơi lội một cách vui vẻ.

``Ối chà, Roboco-senpai cũng có thể bơi được sao?'' Ayame hỏi trong sự ngạc nhiên. Nàng
Người máy đáp: ``Chất này không đủ ướt để làm hỏng chị đâu\~{} Với cả bộ trang phục Năm
Mới này cũng đã bảo vệ được chị rồi\~{}''

``Tức là bình thường em sẽ không bơi được chứ gì?'' Kuon hỏi với một giọng đậm chất cà
khịa. Roboco đáp, vung tay vung chân cố biện minh: ``T-Thì em là rô-bốt mà! Làm sao chịu được nước!''

``Câu sau đá câu trước rồi kìa,'' nàng Pháp sư nói, trong khi cố lôi nàng Vu nữ lên. Ayame
thì vẫn mỉm cười tinh nghịch, quay tất cả mọi thứ, bắt trọn mọi khoảnh khắc.

Sau một hồi loay hoay, nhóm Fubuki, Noel, Roboco, và Haato cuối cùng cũng thoát khỏi
chiếc xe đầy bùn đất và bước lên sân nhà trong tình trạng lấm lem, nhếch nhác. Rushia
thì đưa được Miko vào trong nhà ở tình trạng bất tỉnh, mắt vẫn đang xoáy sau cơn sốc vừa
rồi.

Kuon, dẫn đầu nhóm, vừa đi vừa lắc đầu: ``Thật là\dots\ Mấy đứa vào rửa ráy đi, ông bà đang
chờ đấy.'' Cả nhóm rộn ràng tiến vào, sẵn sàng cho buổi tham quan nhà của ông bà nội
vào buổi chiều ngày mùng Một đầy hỗn loạn nhưng đáng nhớ.

\ct

Căn nhà của ông bà nội Kuon, một gian nhà ba tầng cũ kỹ nhưng đầy hoài niệm, mang nét
đặc trưng của những ngôi nhà thời Pháp thuộc xưa, nhưng cũng mang hơi hướng đôi chút
của vùng quê những năm đầu 2000. Bề ngoài được sơn một màu hồng đỏ, với sân phía
trước được lát gạch đỏ, cùng với những chậu cây xanh vẫn còn khá tươi tốt cho chính tay
ông bà trồng. Căn nhà có một cánh cửa đỏ làm bằng thép, tường dựng bằng bê tông và
những gai hoa lá.

Tầng một là nơi sinh hoạt chính, gồm phòng thờ trang nghiêm và hai chiếc ghế cùng một
bàn uống nước chè; phòng khách với vài ba kệ tủ kính, một chiếc giường gỗ không có
chiếu lót; phòng ngủ với một chiếc giường gỗ được rải chiếu cùng một vài vật dụng lặt
vặt khác; và một nhà vệ sinh đơn sơ, bẩn thỉu do chưa dọn, chỉ đủ vừa cho một người.
Bên ngoài, khu vườn nhỏ ngày trước với chuồng gà và chuồng lợn bị cây dại xâm lấn, tạo
nên vẻ hoang sơ đến lạ.

Điều đặc biệt là lối lên tầng hai đã bị chặn bởi một chiếc tủ đồ lớn từ bao năm nay. Gia
đình của nhà nội Kuon, đã quen với cuộc sống ở tầng một kể từ lúc họ trưởng thành, cũng
không bận tâm tìm hiểu tầng trên có gì. Các thành viên tới đây lúc đầu tới đây cũng không
để ý gì lắm tới chúng.

Thế nhưng, nàng Quỷ sừng Ayame thì không nghĩ vậy. Đứng trước chiếc tủ cũ, cô tò mò
quay sang nhóm bạn: ``Nè, nè! Mọi người có muốn xem trên tầng có gì không?''

Nàng Y tá đang ung dung ngồi bên tách trà, lắc đầu nhẹ: ``Cô không nghĩ là trên đó có gì
đáng xem đâu, Ayame-sama.'' Ngay cả cô nàng nhí nhảnh như Haato - người vừa mới lau
xong bùn đất trên người, lên tiếng: ``Với cả đây là nhà người lạ mà. Ta không thể tùy tiện
lên trên đó, đúng chứ?''

Mio dù vẫn còn chút chóng mặt từ ``cuộc đua'', cũng đồng tình: ``Tôi phải đồng ý với
Haato-chan ở khoản này.'' Dẫu cho mọi người nói vậy, nàng Oni không chịu gì bỏ cuộc.
Đôi mắt đỏ long lanh như thể nhìn thấy điều gì đó hấp dẫn: ``Biết đâu trên đó có kho báu
thì sao!?''

Kuon khoanh tay, khẽ lắc đầu: ``Không, không. Anh không nghĩ là gia đình anh có cái gì
quý giá làm kho báu đâu.'' Dù vậy, sự phản đối của mọi người không làm nàng Quỷ sừng nản chí. Nhưng trước khi Ayame kịp làm gì, Roboco lặng lẽ đứng dậy.

Nghe lỏm toàn bộ cuộc trò chuyện, nàng Người máy tiến tới chiếc tủ và\dots\ chỉ mất vài
phút để dọn sạch mọi chướng ngại. ``Chị kê đồ xuống rồi đây\~{}''
Câu nói nhẹ tênh của nàng Người máy khiến mọi người sửng sốt. Trong khi các thành
viên còn chưa kịp phản ứng, Ayame đã mắt sáng hơn cả ngọc, nhảy lon ton lên cầu thang,
bỏ mặc những tiếng gọi ngăn cản từ phía sau.

\begin{dialogue}
  \speak{Kuon} N-Này! Xuống đi! Trên đó không có gì đâu!
  \speak{Ayame} Em phải kiểm tra thử xem trên này có gì không!
  \speak{Mio} Mà này, Roboco-senpai, sao chị lại tự tiện dỡ đồ xuống vậy hả?
\end{dialogue}



Trong lúc Roboco gãi đầu cười khì, Kuon chỉ có thể nhìn nàng Quỷ sừng tinh nghịch kia
đang thoăn thoắt chạy lên trên tầng. Cậu bất lực, xoa trán: ``Thôi xong\dots\ Anh đoán tụi
mình lại phải đuổi theo, đúng không?''

Choco thở dài: ``Nếu không đi thì làm sao biết được Ayame-sama sẽ làm gì trên đó?''
Nàng Sói đen gật đầu: ``Cậu ta là chúa nghịch mà\dots''

Cả nhóm cuối cùng cũng miễn cưỡng theo chân Ayame, bước lên tầng hai với tâm trạng
vừa tò mò vừa lo lắng.

\ct

Tầng hai, dù đã lâu không được sử dụng, vẫn toát lên vẻ cũ kỹ nhưng đầy bí ẩn. Khu vực
này cũ kỹ tới mức, phần sơn đã bị bong tróc hoàn toàn, chỉ còn lại phần bê tông cứng cáp.
Nơi đây gồm năm phòng, hai trong số đó bị khóa chặt, chỉ có thể nhìn vào bên trong qua
cửa kính ra vào khá lớn, và một thì phòng cũng khóa chặt, nhưng có thể luồn tay vào bên
trong.

Cánh cửa phòng thờ là một trong hai phòng mà Ayame có thể mở ra, và khi cô tò mò ngó
vào, cô nhìn thấy một bàn thờ lớn trang nghiêm được đặt ngay cạnh cửa. Chiếc bàn thờ
gỗ này về cơ bản giống chiếc ở phòng thờ tầng một, còn về bản thân căn phòng thì nó
cũng chẳng có gì cả.

Ayame vội đóng cửa lại, lẩm bẩm: ``Ể\~{} Chẳng có gì cả.'' Đúng lúc đó, mọi người cũng
vừa chạy lên đến nơi. Kuon, vẫn còn thở dốc sau khi đuổi theo nàng Oni, cất giọng gắt
gỏng: ``Đã bảo rồi, trên này không có gì đặc biệt đâu!''

Mio thì nhẹ giọng hơn: ``Tôi nghĩ chúng ta nên xuống đi\dots'' Tất cả những người còn lại
cũng gật đầu.

``Chị thấy nơi này xám xịt thế nào ấy\dots'' Haato khẽ rùng mình khi nhìn xung quanh hành
lang của tầng hai này, khi ngoài ô cửa sổ lớn ở phía cầu thang chiếu vào và từ các cửa
phòng, còn lại nơi này cũng không có đèn đóm gì cả.

Tuy nhiên, nàng Quỷ sừng dường như không chịu từ bỏ. Cô cố thử mở ba căn phòng còn
lại, nhưng đều không thành công.

``Sao\dots\ tôi\dots\ không mở được\dots?'' Ayame gồng hết sức để mở chúng, nhưng không tài
nào mở nổi. Choco phỏng đoán: ``Cô nghĩ chắc là người nhà ở đây đã khóa chúng đi do
không sử dụng rồi.''

``Với cả, kể cả mở nó ra thì cũng không có gì trong này mà,'' nàng Sói đen thêm thắt. Qua
cửa kính, họ thấy bên trong hai căn phòng trống rỗng, chỉ còn sót lại khung giường cũ.

Phòng còn lại thì chất đầy đồ đạc cũ kỹ, bụi phủ mờ mịt, có vẻ như đã bị biến thành nhà
kho từ rất lâu. Tưởng chừng đống đồ đó chỉ coi như là rác với đa số mọi người, nhưng
với Ayame\dots

``Trong này có thể có đồ mà ta sẽ dùng tới này!''

\dots mắt của cô bỗng chưng sáng lên, như thể cô vừa tìm thấy kho báu vậy. Haato đứng nép
sang một bên, lắc đầu đầy lo lắng: ``Chị không nghĩ đây là ý hay đâu\dots''

Cậu Tổng thì chỉ thở dài: ``Trông em chẳng khác gì một đứa trẻ lên mẫu giáo vậy.''

Bỏ ngoài tai những lời khuyên can, nàng Oni nhanh nhẹn trèo qua cánh cửa nhỏ, nhảy
vào căn phòng chất đầy đồ đạc, lục lọi như đang tìm kho báu. ``Để xem trong này có đồ
gì hữu ích cho chúng ta nào!'' Cả bốn người nhìn cô nàng đang hiếu kỳ kìa, thở dài bất
lực.

``Nhân tiện, đây cũng là phòng của bác gái trưởng nhà anh đấy,'' cậu thuận tiện giới thiệu
cho các cô gái.

``Vậy căn phòng bên cạnh này?'' nàng Song tính chỉ tay về phía một căn phòng khác.

``Phòng bố anh ngày xưa.''

``Còn căn phòng đối diện với chúng ta đang đứng này?'' Mio thắc mắc?

``Phòng bác gái thứ.''

Roboco quay lại nhìn phòng thờ: ``Còn đây? Ayame vẫn đang lục lọi trong phòng kho,
ngẩng đầu lên: ``Em nghĩ đó là phòng thờ đấy. A, có lò xo này!''

Nàng Quỷ sừng giơ ra một cuộn lò xo còn sử dụng tốt, đôi mắt sáng rỡ như vừa phát hiện
ra kho báu. Tuy nhiên, khi nhìn thấy món đồ này, cả nhóm chỉ biết ngẩn ra nhìn cô, rồi
nhìn nhau, ánh mắt như muốn hỏi: ``\textit{Cái đó thì làm được gì?}''

``Vậy còn căn phòng trông ẩm ướt đằng kia\dots?'' nàng Y tá chỉ về căn phòng nơi bức tường
được lát gạch gốm xanh lam toát lên vẻ hiện đại, nhưng phần sàn thì vương vấn bụi bẩn.

``Nhà vệ sinh. À mà, hôm nọ các bác ấy về dọn qua nhà, có phát hiện một con chuột chết
trong đó đấy,'' cậu nói với mọi người với vẻ thản nhiên. Cả nàng Sói và nàng Y tá tái xanh
mặt mày, còn nàng Quỷ sừng, nàng Song tính và nàng Người máy thì cũng chỉ ``ồ'' một
tiếng.

``Đừng lo, các bác ấy vứt đi rồi.'' Câu nói của cô khiến Choco và Mio thở phào nhẹ nhõm.
Dẫu vậy, Sói đen vẫn run rẩy: ``E-Em không muốn vào trong đó đâu\dots''

``Đồng ý.''

\ct

Sau khi giới thiệu qua tầng hai, Kuon tiện thể kể thêm về tầng ba: ``Tầng trên có hai
phòng, một là của bác Hijaki, và một là phòng dẫn ra ban công nhìn xuống vườn.''

Cả nhóm tiếp tục trèo lên tầng ba, không quên nhìn quanh và đặt ra nhiều câu hỏi. Mio
đứng tựa vào lan can, phóng mắt nhìn ra xa, khẽ cảm thán: ``Không khí ở đây thật dễ chịu.
Khác hẳn trong thành phố\dots''

Kuon khẽ cười: ``Vì đây là thôn quê mà. Ở quê dù gì cũng dễ thở hơn so với ở trung tâm
thành phố.'' Bất chợt, Haato hỏi: ``Mà, nhà anh trông đông người như thế, ngủ nghỉ sao
cho đủ nhỉ?''

Cậu con trai nở nụ cười nhẹ, giải thích rõ hơn: ``Thực chất ra đây là nhà của ông bà nội
anh, và anh chỉ chơi loanh quanh ở tầng một thôi. Nhưng mà anh nhớ bố anh có kể thế
này\dots''

Cậu tiếp tục giảng giải về lịch sử của ngôi nhà tại thôn quê này. Theo lời của bố cậu kể,
gia đình Hikawa nếu tính ông bà nội cậu là già nhất, thì nhà cậu có năm người con: ba
trai và hai gái. Bố cậu là con trai út, Hijaki là con trai thứ.

``Vậy, bác trai trưởng nhà anh sẽ ngủ ở đâu?'' Roboco thắc mắc?

Kuon từ tốn đáp, với một tông giọng nhẹ nhàng hiếm hoi: ``À, anh cả trong nhà thì được
ông bà ưu tiên ngủ dới tầng một. Mấy đứa đã xem qua cái phòng ngủ ở dưới đó rồi chứ?''

Nàng Y tá tỏ vẻ ngạc nhiên: ``À, vậy tầng một cũng là nơi bác trưởng nghỉ ngơi sao,
Kuon-sama? Vậy, hai bác gái thì sao?''

Cậu chỉ tay về phía tầng hai, nơi họ vừa xem qua: ``Hai phòng ở tầng hai, một là của bác
gái trưởng, và một là của bác gái thứ. Họ không còn ở đây, nhưng vẫn để lại đồ đạc cũ.
Như phòng mà Ayame-chan đang lục lọi đấy, đó là phòng của bác gái trưởng.''

Nàng Oni từ bên trong căn phòng đầy đồ ló đầu ra, tay vẫn cầm cuộn lò xo, nói đầy hứng
thú: ``Bác ấy mà về đây, chắc sẽ ngạc nhiên lắm vì ta tìm thấy kho báu của bà ấy đấy, yo!''

Mọi người chỉ biết bật cười, còn Haato thì quay sang Kuon tiếp tục tò mò: ``Thế còn ông
bà nội anh? Họ ở đâu?''

Kuon nhún vai: ``Họ ở phòng dưới tầng một. Họ dù gì cũng lớn tuổi nhất mà, nên cần
một nơi tiện nghi và dễ đi lại hơn.''

Nàng Người máy gật gù, nhìn quanh tầng ba: ``Vậy tầng này chỉ có phòng của bác Hijakisan
thôi ạ?'' Kuon khẽ gật đầu: ``Ừ. Căn phòng đó giờ cũng chẳng còn gì đâu. Bác ấy bây
giờ đang ở với bọn anh đấy, chắc Mio-chan cũng đã gặp rồi nhỉ?''

``Em nghĩ là vậy\dots'' nàng Sói đen nói, tỏ vẻ thắc mắc. Khi mọi người ngó sang căn phòng
đã từng là nơi ở của bác trai thứ nhà Kuon, nó đã sờn màu do thời gian, và chỉ còn lại
khung giường còn ở đó.

Cậu chỉ tay về phía cánh cửa đối diện: ``Còn căn ban công kia\dots\ anh nghĩ là thường chẳng
ai lên cả. Anh đoán chắc đó là nơi họ phơi quần áo chăng? Thỉnh thoảng lắm bọn anh mới lên ngắm nghía rồi đi dọn tí xíu.''

Nàng Oni nghe thấy vậy, lập tức lấp ló ra khỏi phòng, hào hứng giơ cuộn lò xo lên: ``Nếu
thế thì từ nay ta nhận nhiệm vụ canh giữ tầng ba nhé!''

Cô lon ton chạy lên trên tầng ba với mọi người, rồi cùng nhau bước vào. Đúng như cậu
nói, xung quanh họ là một khu khá rộng với bàn ghế đã mai mọt đi, xếp một cách lộn
xộn kính thì đầy vết bẩn, sàn thì còn vương vất vài chiếc lá rơi, cốc thì cũng bẩn như thế,
nhưng chưa bị lên mốc.

Còn lại, khung cảnh trước mặt họ là một khu vườn được trải rộng ra, với hàng tá cây dại
mọc lên um sùm cùng vài ba cái cây cao mọc lạc tại đây. Ngoài một quang cảnh ngoài
vườn, họ còn có thể nhìn thấy khung cảnh làng quê thanh bình, nơi không có nhà cửa cao
tầng, chỉ có thấp thoáng vài xe công nông và vẫn còn vài người đi bộ quanh khu làng.

Thấp thoáng bên ngoài, Mel và Rushia đang đi dạo quanh khu vườn nhà. Một người bước
đi nhẹ nhàng với dáng vẻ như đang thưởng thức không khí trong lành, người còn lại thi
thoảng dừng lại ngắm nhìn những đóa hoa dại mọc lưa thưa bên vệ đường. Dáng vẻ lặng
lẽ nhưng dễ thương của cả hai khiến nhóm người đứng trên tầng ba không khỏi mỉm cười.

``Cảnh này đẹp thật đấy. Chẳng mấy khi ta được thấy một nơi bình yên như thế này\dots''
Ayame lẩm bẩm. Choco đứng cạnh nàng Quỷ sừng, cũng gật gù đồng tình: ``Nhìn Melsama
và Rushia-sama kìa. Đúng là không khí thôn quê có thể làm con người ta thư thái
hơn hẳn.''

Haato chống tay lên lan can, nhìn xa xăm rồi cất giọng: ``Nếu chúng ta ở đây thêm vài
ngày nữa, chắc cũng quên hết mấy chuyện xui xẻo hồi sáng mất thôi.'' Tất cả mọi người
bật cười trước câu nói có phần ngây thơ của cô nàng.

Mio đưa mắt ngắm nhìn bầu trời trong xanh không chút gợn mày, khẽ nói như tự nhủ:
``Đúng vậy\dots\ Không khí ở đây làm ta cảm thấy mình đang được trẻ lại vậy.''

Kuon đứng tựa vào cánh cửa, khẽ thở dài: ``Anh cũng ước cho mấy đứa ở lại đây vài ngày
lắm, nhưng mà mấy đứa sau hai ngày nữa vẫn phải đi làm rồi nhỉ. Với cả, anh không chắc
các em chịu được việc không có mạng để stream đâu.''

Nàng Người máy bật cười: ``Anh không biết là Ojou-san đang định biến nơi này thành
`căn cứ bí mật' sao? Chắc cô ấy chẳng cần mạng đâu.'' Ayame đứng phắt dậy, chỉ tay ra
phía làng quê: ``Đúng vậy! Từ nay nơi này sẽ là pháo đài của Ayame-sama! Mọi người
nhớ báo cáo trước khi vào đây nha, yo!''

Không ai nhịn nổi cười trước sự hùng hồn của nàng Quỷ sừng. Bầu không khí vui vẻ lại
một lần nữa tràn ngập trên ban công tầng ba.

\ct

Ở dưới vườn, Noel vẫn đang loay hoay chặt cây dại với vẻ mặt đầy quyết tâm. Cô lẩm
bẩm trong lúc quơ lưỡi dao phát mạnh: ``Nhà này nhiều bụi cây quá\dots\ Không biết lần
cuối cùng ai dọn là khi nào nữa.''

Choco xuống dưới nhà, nhìn nàng Hiệp sĩ đang miệt mài chặt bỏ một trong những cây cao kia bằng một con dao bầu mà cô tìm được trong nhà. Cô đáp: ``Cũng phải. Lâu rồi
gia đình anh ấy chưa dọn mà. Mà em có chắc là mình cần chặt hết không đấy?''

Noel, tạm dừng tay để lau mồ hôi, nhún vai: ``Em không nghĩ là với sức trâu của mình thì
em có thể bỏ hết được đâu\dots\ Không khéo em phải đốt cả khu vườn này mất thôi.''

Nàng Y tá giật mình: ``Đừng. Nhỡ đâu em đốt luôn cả căn nhà của anh ấy thì sao?'' Nàng
Hiệp sĩ nhìn quanh một lượt, chỉ tay về phía một khu vực nhỏ vốn dĩ là bếp ngoài trời:
``Nhìn xem, khu vực này trông có vẻ chẳng còn ai dùng đến. Em thấy đốt cũng không ảnh
hưởng gì.''

Khu bếp đó đúng là đã tan hoang, chỉ còn lại vài chạn bát rỉ sét và một chiếc máng rửa
rau quả phủ bụi, như thể nó đã bị lãng quên qua nhiều thế hệ.

Choco thở dài, nhìn căn bếp đã bị bỏ hoang đó: ``Ừ thì, đúng là nó chẳng còn dùng được
nữa, nhưng đây đâu phải là nhà của chúng ta. Em làm vậy lỡ gây phiền phức thì sao?''

Noel nghe vậy, hơi cúi đầu và nhăn mặt than vãn: ``Thế giờ em phải làm sao? Chặt mãi
thế này cũng không ổn.'' Nhìn cô nàng đang thở ngắn thở dài như vậy, Choco chống tay
lên cằm: ``Cô nghĩ là\dots''

Câu nói của nàng Y tá còn chưa dứt thì từ sâu trong khu vườn, bỗng có tiếng *rắc rắc*
vang lên. Một số cây lớn bị bật gốc, mật cây bắn tung tóe, rồi chúng từ từ đổ ầm xuống
bên phải, để lộ những chiếc rễ lớn chằng chịt bên dưới.

Cả Noel và Choco đều sửng sốt, mắt mở to nhìn cảnh tượng trước mặt. Tiếng động lớn
này thậm chí còn thu hút sự chú ý của nhóm người trên tầng ba. Mio, nhìn xuống từ ban
công, thốt lên: ``Gì thế kia!? Có chuyện gì xảy ra vậy!?''

Nàng Hiệp sĩ sau vài giây sững sờ, không hiểu chuyện gì vừa xảy ra, đột nhiên tươi tỉnh
hẳn: ``May quá, em không phải làm nữa!''

Trong khi đó, nàng Y tá vẫn lẩm bẩm, không khỏi ngỡ ngàng sau sự việc kỳ lạ đó: ``Chuyện
quái gì vừa xảy ra thế nhỉ\dots''

Lúc này, Sora từ trong nhà bước ra, nhẹ nhàng lên tiếng như đang giải thích: ``Em nghĩ
chắc là bởi\dots\ ngực của hai người đấy.''

Câu nói của nàng Thần tượng khiến cả Noel lẫn Choco sững sờ, đồng loạt đỏ mặt. Họ
nhìn nhau, vừa ngượng ngùng vừa không hiểu vì sao chuyện kỳ lạ này lại liên quan đến\dots
vòng một của mình.

Ở tầng ba, nhóm người quan sát cũng không nhịn được cười. Haato ngã người dựa vào
lan can, cười lớn: ``Sora-senpai đúng là có cách lý giải độc đáo thật đấy!''

Không khí lại trở nên sôi động hơn bao giờ hết, trong khi Noel và Choco cố gắng lấy lại
bình tĩnh, tiếp tục ``cuộc chiến'' dọn vườn.

\ct

Khung cảnh làng quê thanh bình với bầu trời xanh trong và những cánh đồng trải dài xa
tít, điểm xuyết vài ngôi nhà nhỏ lợp mái ngói đỏ khiến không gian tràn ngập vẻ yên bình mà hiếm khi họ tìm thấy ở thành phố.

Pekora, Miko, Fubuki, Matsuri, Okayu và Korone hứng thú tản bộ quanh làng, tò mò quan
sát mọi thứ tại nơi mà họ hầu như chưa bao giờ tới.

Nàng Thỏ nhấp nhô đôi tai khi cô bước đi, nhí nhảnh chỉ tay về những cánh đồng xa xa:
``Peko\~{}! Chỗ đó có gì thế nhỉ? Có phải ruộng cà rốt không?''

Nàng Vu nữ vừa cười vừa nhún vai: ``Không đâu, Peko-chan, ne\~{} Đấy là lúa mà, làm gì
có cà rốt nào mà mọc cao như thế?''

Nàng Cáo trắng đang chụp ảnh bằng điện thoại, quay sang góp chuyện: ``Nhưng chắc
chắn là nhìn đẹp thật. Chúng ta nên chụp vài tấm ảnh làm kỷ niệm!''

Trong khi đó, Matsuri và Okayu mải mê khám phá khu chợ nhỏ gần đó. Korone thì chạy
tung tăng bên cạnh, không ngừng vẫy tay chào những người dân làng thân thiện.

Dẫu cho họ gặp một chút khó khăn vì ``giao diện'' có phần khác lạ, người dân tại khu làng
này cũng không mảy may bận tâm tới, đoán rằng họ cũng chỉ là những cô gái thích hóa
trang.

\ct

Ở nhà, Sora và Mel ở lại để chăm sóc ông bà nội của Kuon. Sora cẩn thận rót trà, trò
chuyện với ông bà, trong khi Mel nhẹ nhàng dọn dẹp góc bàn thờ.

Nàng Thần tượng mỉm cười dịu dàng hỏi: ``Ông bà có cần chúng cháu giúp gì nữa không
ạ?'' Hai cụ già chỉ mỉm cười lắc đầu, tỏ vẻ hài lòng khi thấy không khí nhộn nhịp và ấm
cúng trong nhà. Nàng Dơi vừa lau chiếc lư đồng sáng bóng, vừa nhìn ra ngoài cửa sổ,
cảm nhận sự an lành hiếm có nơi đây.

\ct

Trong khi đó, Aqua, Shion và Subaru tiếp tục công việc lau dọn nhà cửa. Nàng Hầu gái
hăng hái báo cáo, với cây chổi trên tay: ``Tôi quét lá xong rồi đây- Ồ\dots''

Cô dừng lại khi nhìn thấy đống bụi và mạng nhện đầy góc nhà. Shion, đang cầm cây chổi,
gật đầu như hiểu ra vấn đề: ``Có vẻ như chúng ta lại phải dọn cả cái nhà này rồi nhỉ.''

Không đợi thêm giây phút nào, nàng Phùng thỷ vung tay niệm một phép thuật nhỏ, giống
như cách cô đã làm khi dọn phòng của Hijaki. Tuy nhiên, Subaru nhận ra ngay nguy cơ:
``K-Khoan đã, Shion-chan! Dừng lại! Bà làm thế thì-''

Nhưng mọi chuyện đã quá muộn. Một làn bụi khổng lồ bốc lên, kèm theo hàng loạt tiếng
*bụp bụp* kỳ lạ. Đống bụi không biến mất mà chỉ\dots\ chuyển từ góc này sang góc khác,
thậm chí phủ thêm cả tầng một.

Aqua ho sặc sụa, vẫy tay trước mặt để xua bụi: ``Shion-chan, bà đang làm kiểu gì thế
này!?'' Shion nhún vai: ``Thì mấy người có chỉ tôi rõ là phải làm thế nào đâu.''

Câu trả lời của cô khiến nàng Vịt ngã ngựa ra sau. Cô thảng thốt: ``Thế mà tôi tưởng bà
tinh tướng như thế nào chứ!'' Cô đứng giữa đám hỗn độn, chỉ biết ôm đầu thở dài: ``Đến cả nơi làng quê cũng lắm chuyện oái oăm nữa!''

Kuon vừa từ tầng ba xuống, nhìn cảnh tượng bừa bộn, chỉ có thể ngán ngẩm lắc đầu:
``Đúng như em nói đấy, Subaru-chan ạ.'' A-chan cầm máy quay gần đó, cũng ngán ngẩm:
``Tôi phải đồng ý với cậu, Hikawa-san\dots\ Đúng là không chỗ nào yên bình với chúng ta
cả\dots ''

Sau khi tạm thời dọn dẹp hậu quả, nhóm còn lại quyết định đi vòng quanh làng. Họ tò mò
tìm hiểu cuộc sống hàng ngày của người dân địa phương, từ việc cấy lúa, chăn trâu cho
đến những món đồ thủ công đơn giản mà đầy ý nghĩa.

Cảm giác thư thái và bình dị ấy như xóa tan mọi mệt mỏi, khiến họ quên đi những phiền
muộn thường ngày ở thành phố.

Tối hôm đó, cả nhà và các cô gái \hll\ cùng nhau chuẩn bị bữa cơm tối. Trong không
khí rộn ràng của ngày đầu năm mới, căn bếp trở thành trung tâm của những tiếng cười
nói, tiếng chặt thịt, tiếng rán nem và cả những câu chuyện không hồi kết.

\ct

Tối hôm đó, cả nhà và các cô gái \hll\ cùng nhau chuẩn bị bữa cơm tối. Trong không
khí rộn ràng của ngày đầu năm mới, căn bếp trở thành trung tâm của những tiếng cười
nói, tiếng chặt thịt, tiếng rán nem và cả những câu chuyện không hồi kết.

Fubuki luôn lanh lợi, đảm nhận việc rửa rau và sơ chế nguyên liệu. Cô vừa làm vừa vui
vẻ trêu đùa: ``Này, có ai thấy cái này giống anime tập Tết không? Cứ như chúng ta đang
quay show vậy!''

``Thì bởi vì đúng là chúng ta đang quay một bộ phim về Tết mà,'' Kuon nói, cùng lúc đang
sắp bát đũa chuẩn bị dọn cơm.

Mio đứng bên cạnh giúp cuốn nem, nhưng đôi tay vụng về khiến một số chiếc nem bị
rách. Subaru thấy thế liền phá lên cười: ``Mio-sha, làm thế nào mà bà còn tệ hơn tôi được
chứ?''

Nàng Sói đen ngượng ngùng đáp trả: ``K-Không phải lỗi của tôi đâu, tại miếng bánh tráng
này khó cuốn quá mà!''

Noel thì chẳng bận tâm đến việc cuốn nem hay nấu nướng. Cô tập trung toàn lực vào
việc\dots\ canh chừng gà luộc, đôi mắt sáng rực mỗi lần nhìn vào nồi nước sôi bốc hơi nghi
ngút. Bụng của cô thỉnh thoảng còn kêu lên inh ỏi vì đói.

``Sắp được ăn rồi, sắp được ăn rồi\dots'' nàng Hiệp sĩ lẩm bẩm, mơ tưởng về món gà luộc
sắp sẵn sàng. Rushia, đứng xa xa, lặng lẽ nhìn mọi người làm việc. Khi thấy Kuon bận
rộn chỉ đạo từng người, nàng nhỏ giọng hỏi: ``Kuon-senpai, có cần em giúp gì không?''

Cậu Tổng quay lại, mỉm cười: ``Không cần đâu, em cứ nghỉ ngơi đi. Nhưng nếu muốn,
em có thể giúp chị A-san thái thịt đông nhé.''

Ở góc bếp khác, Matsuri và Okayu đang tranh nhau ai sẽ rán nem trước.

\begin{dialogue}
  \speak{Matsuri} Chị muốn được rán nem trước!
  \speak{Okayu} Nhưng mà em cũng muốn rán mà\~{}
  \speak{Matsuri} Nhưng mà Ryuuhou-chan nói nem của chị rán ngon hơn!
  \speak{Okayu} Đâu, em mà\~{}
\end{dialogue}

Korone, như thường lệ, hào hứng cổ vũ: ``Để cả hai làm cùng luôn đi! Nem của hai người
chắc chắn sẽ ngon nhất!''

Aki, với dáng vẻ dịu dàng, đảm nhận phần trang trí xôi gấc lên đĩa.

``Xôi gấc này mình nên để đâu nhỉ\dots'' nàng Dị giới lẩm bẩm. Cô quay sang cậu Tổng,
hỏi: ``Này, Kuon-senpai, em nên để nó ở đâu?''

``Ở giữa bàn ấy. Tiện nhất thì chỗ này nhét vừa thì để,'' cậu đáp lại.

Choco thì thử nêm nếm món thịt đông, nhưng lại làm ra vẻ ``đầu bếp chuyên nghiệp'' khiến
mọi người cười ồ: ``Món thịt đông này đúng là ngon tuyệt vời!''

Kuon đổ mồ hôi hột: ``Cái món đó là mẹ anh làm mà, có phải là em làm đâu\dots''

Roboco đứng cạnh cậu, thở dài: ``Choco-sensei, cô đang nếm hay là ăn luôn đấy? Món
này phải để cả nhà cùng thưởng thức mà.''

Haato thì đang bận\dots\ tự chiến đấu với một chiếc bánh chưng khó cắt. ``Cái bánh chưng
này sao mà dai ngoách vậy!?'' A-chan thấy vậy liền giúp, vừa làm vừa lắc đầu cười:
``Akai-san, cô dùng dao sai rồi đấy. Để tôi chỉ cho cô cách.''

Trong khi đó, Mel đang trò chuyện với Sora ở góc bàn gần bếp. Nàng Dơi nhìn mọi
người cùng nhau nấu ăn, nở một nụ cười mãn nguyện: ``Không khí ở đây thật khác hẳn
với Tokyo nhỉ, Sora-senpai?''

Nàng Thần tượng mỉm cười nhẹ nhàng: ``Ừ, bình yên và giản dị hơn nhiều. Chị thấy rất
dễ chịu.''

Pekora, người đang phụ cắt rau củ, bất chợt ngẩng đầu lên với một nụ cười lém lỉnh:
``Kuon-senpai, anh có định làm món gì đặc biệt kiểu `bí mật gia truyền' không đấy, peko?''

Cậu con trai không quay đầu lại, chỉ đáp ngắn gọn: ``Không có gì bí mật đâu, nhưng em
cứ chờ xem.''

Shion ở gần đó, liền xen vào: ``Món ăn gia truyền mà không bí mật thì mất hay rồi. Hay
để em thêm một chút `ma thuật' vào cho đặc biệt hơn nhé?''

Aqua đang quét dọn ở góc bếp, vội kêu lên: ``S-Shion-chan, đừng có làm gì kỳ quặc nữa!
Lần trước bà đã biến cả nồi súp tôi làm thành màu tím hết rồi đấy!'' Đáp lại, nàng Phù
thủy vờ như không nghe thấy gì, chỉ mỉm cười đầy ẩn ý.

\ct

Khi bữa tối bắt đầu, bàn ăn trở nên nhộn nhịp với các món đặc sản ngày Tết của vùng
Etsunan. Noel là người nổi bật nhất, với tốc độ ăn gà luộc nhanh đến mức làm mọi người
tròn mắt.

Nàng Vịt cười lớn, giở giọng trêu: ``Noel-chan! Em ăn như thể chưa được ăn gà bao giờ
ấy!'' Nàng Hiệp sĩ không đáp, chỉ ăn món thịt gà đầy ngon lành, như thể đã làm tâm
nguyện của cô lâu năm được hoàn thiện.

Mio cắn một miếng nem rán do Matsuri và Okayu làm (được Ryuuhou hướng dẫn), gật gù
đồng tình: ``Ryuuhou-chan nêm nếm vừa miệng thật. Chị phải học cách làm mới được!''

Một nửa số thành viên thích thú khám phá món bánh chưng, trong khi nửa còn lại thì hơi
chần chừ với hương vị lạ lẫm. Korone, với tinh thần ``không ngán món gì'', mạnh dạn ăn
một miếng lớn rồi thốt lên: ``Ngon quá! Mọi người thử đi, đừng sợ!''

Nàng Mèo bên cạnh bật cười: ``Koro-san làm cứ như thể mọi người vẫn còn lạ lẫm với món đó
á\~{}''

Ở góc bàn, Sora và Mel đang gắp thử món thịt đông. Nàng Dơi vừa ăn vừa nhíu mày tò
mò: ``Món này lạ thật, nhưng cũng thú vị.'' Nàng Thần tượng mỉm cười: ``Ừ, có chút mát
lạnh, ăn rất hợp vào dịp Tết.''

Bữa ăn kết thúc trong sự thỏa mãn của tất cả mọi người. Không chỉ vì các món ăn ngon,
mà còn bởi bầu không khí gia đình đầm ấm và những tiếng cười không ngớt.

Mùng Một Tết kết thúc với bao sự hỗn loạn đến ngớ ngẩn và phi logic, nhưng cũng đầy ý
nghĩa. Ngày Tết, dù ở bất cứ đâu, vẫn luôn mang đến cảm giác sum vầy và ấm áp khó tả.

\ct

\pov{Hikawa Kuon}

\dots

\dots

\dots Hửm?

Tôi bất giác bừng tỉnh giữa đêm, cảm nhận rõ sự chật chội bất thường quanh mình. Đèn
ngủ bên góc phòng tỏa ánh sáng dịu nhẹ, nhưng đủ để tôi nhận ra khung cảnh trước mắt.

Bảy người con gái, mỗi người một tư thế, đang nằm trên tấm nệm chật chội, tất cả đều
ngủ say sưa như không hề hay biết điều gì.

Miko-chan, nằm co tròn như mèo, đầu gối sát vào cánh tay của tôi. Noel-chan thì chiếm
phần lớn diện tích giường bằng thân hình khỏe mạnh, tay vắt qua vai như một cử chỉ vô
thức. Thậm chí trông như thể còn đá tôi ra khỏi đệm nữa. Rushia-chan nhẹ nhàng nhất,
nằm gọn gàng bên cạnh tôi, gần như không động đậy.

Roboco-chan nằm vắt ngang chân tôi, với dáng vẻ vô tư như một đứa trẻ. Pekora-chan,
dù đang ngủ, vẫn giữ nụ cười nghịch ngợm trên môi, tay quàng lấy chiếc gối mà chúng
tôi có. Okayu-chan thì ôm lấy một con gấu bông nhỏ, mặc dù em ấy hoàn toàn tựa lưng
vào tôi. Và cuối cùng, Korone-chan, với vẻ mặt hiền lành, nằm úp bụng xuống giường,
mái tóc xõa tung che nửa khuôn mặt.

Tôi nhìn bảy cô nàng ngủ một cách bình yên đó, khẽ thở dài: ``\textit{Lại là mình sao\dots}'' Lời nói
của tôi rất nhỏ nhẹ, cố gắng không đánh thức ai.

Từ phía bên ngoài căn phòng chật chội ở nhà Kyuukami, ánh đèn ngoài cửa hắt vào khung
cảnh của văn phòng, nơi mọi người đã trở về sau ngày mùng Một Tết Cha đầy kỷ niệm.
Không khí vẫn còn phảng phất mùi gà luộc và nem rán từ bữa tối, pha lẫn hương trầm nhè
nhẹ từ phòng thờ mà Sora-chan và Mel-chan đã lau dọn.

``\textit{Chẳng lẽ đây là phần thưởng của cuộc thi làm bánh mochi mà họ mong muốn dây sao?}''
tôi tự nhủ, rồi bỗng khựng lại khi nhớ lại lời hứa của mình.

Quả đúng là như vậy, tôi thực sự đã đồng ý để bảy người chiến thắng ngủ lại cùng một
đêm\dots

Khẽ điều chỉnh tư thế, tôi nhắm mắt lại, cố gắng tận hưởng chút yên bình hiếm hoi. Ít ra,
ngày mai là mùng Hai Tết Mẹ. Mong rằng nó sẽ đỡ hỗn loạn hơn\dots

\dots Nhưng các bạn biết đấy, với \hll, điều này dường như là bất khả thi. Chắc tôi sẽ
phải sẵn sàng đối đầu với những trò đùa nhây của họ thôi\dots


\end{document}